\documentclass[../Main.tex]{subfiles}


\begin{document}


\chapter{Probability and Measure Theory}

Probability theory provides a mathematical framework for studying random phenomena. Many processes in nature, science, engineering, and finance involve outcomes that cannot be predicted with certainty. Examples include the outcome of a coin toss, the arrival time of customers in a queue, or the evolution of asset prices in financial markets. Probability theory allows us to model such situations and to reason quantitatively about uncertainty.

The modern formulation of probability theory is based on \textbf{measure theory}. In measure theory, the goal is to assign a non-negative number to certain subsets of a given set. This number is called a \textbf{measure} and represents a generalized notion of size, extending familiar ideas such as length, area, or volume. A fundamental property of a measure is \textbf{countable additivity}: if a collection of sets is disjoint, the measure of their union is the sum of their measures.

For mathematical reasons, it is not possible in general to assign a measure to every subset of a set. Instead, one restricts attention to a special collection of subsets called a \sfield. A set equipped with a \sfield is called a \textbf{measurable space}, and measures are defined only on these measurable sets.

Probability theory can be viewed as a special case of measure theory, where we are assigning a measure that represents our belief of occurrence of events. This measure is called a \textbf{Probability measure} and it has the additional property that the total measure of the entire space is equal to one. The triplet $(\probspace)$ is called a \textbf{probability space} where : 

\begin{itemize}
    \item The set $\Omega$ is the \textbf{sample space}. It contains all possible outcomes of the experiment. Once an outcome $\omega \in \Omega$ is fixed, all random quantities associated with the experiment are completely determined.
    \item The collection $\A$ is a \sfield of subsets of $\Omega$, called \textbf{events}. These are the sets to which probabilities can be assigned. From the point of view of measure theory, they are precisely the measurable sets. An event $A \in \A$ represents the occurrence of a given property of the outcome $\omega$.
    \item The function $\PP : \A \to [0,1]$ is a \textbf{probability measure}. For each event $A \in \A$, the value $\PP(A)$ represents the probability that the event $A$ occurs.
\end{itemize}

This axiomatic formulation of probability theory was introduced by \textbf{Andrey Kolmogorov in 1933}. The Kolmogorov axioms provide a simple and powerful foundation for probability and allow many results to be developed in a unified and rigorous way.

Although probability theory relies on measure theory, the goals of the two subjects are different. Measure theory studies general notions of size and integration, while probability theory focuses on modeling and analyzing randomness. In this chapter, we mainly develop probability theory while referring to measure-theoretic concepts whenever they are needed.

A prior knowledge of measure theory can be helpful but is not strictly required. The reader is assumed to have seen basic notions such as \sfields, measurable functions, and measures. Whenever these ideas are used, we briefly recall their meaning and explain their role in probability.

The aim of this chapter is to present the fundamental concepts of probability theory in a clear and structured way, combining mathematical rigor with intuitive interpretations based on random experiments.

\medskip
\textbf{Important remark.}
In many situations, the explicit form of the probability measure $\PP$ will not be specified. Instead, we focus on the properties of measurable functions defined on $(\measspace)$, which are called \emph{random variables}.

\medskip
\textbf{Terminology.}
As in measure theory, sets of measure zero play an important role in probability theory.  A property depending on $\omega \in \Omega$ is said to hold \emph{almost surely} (a.s.) if it holds for all $\omega$ in an event of probability one. Thus, \textit{almost surely} has the same meaning as "almost everywhere" in measure theory.

\medskip
\textbf{Notation.}
For a random variable $\bX$, the notation $\PP(\bX \in A)$ means :
\[ \PP(\{\omega \in \Omega : \bX(\omega) \in A\}). \]
\newpage


\section{General Definitions}
In this section, we introduce the fundamental objects of probability theory. We begin with the definition of a probability space, which provides the mathematical framework for modeling random experiments. We then define random variables, which allow us to describe and analyze numerical quantities generated by randomness.

These definitions form the basis of all subsequent results in this chapter. They are stated in a general and rigorous setting, with close connections to measure theory whenever necessary.



\subsection{Probability space}
\defnr{Probability space }{probability_space}{
A \textbf{probability space} is a triple $(\probspace)$, where:
\begin{itemize}
    \item $\Omega$ is a non-empty set called the \textbf{sample space},
    \item $\A$ is a \sfield of subsets of $\Omega$, and
    \item $\PP: \A \to [0,1]$ is a \textbf{probability measure}.
\end{itemize}
}
We remind the reader that a \sfield $\A$ is a collection of subsets of $\Omega$ such that:
\begin{itemize}
    \item $\Omega \in \A$,
    \item if $A \in \A$ then $A^c = \Omega \setminus A \in \A$,
    \item if $A_1, A_2, \ldots \in \A$ then $\cup_{i=1}^\infty A_i \in \A$.
\end{itemize}

\defnr{Probability measure }{probability_measure}{
A \textbf{probability measure} $\PP$ on a \sfield $\A$ of a set $\Omega$ is a function $\PP: \A \to [0,1]$ satisfying:
\begin{itemize}
    \item $\PP(\Omega) = 1$
    \item For any countable sequence of disjoint sets $A_1, A_2, \ldots \in \A$,
    \[
    \PP\left( \bigcup_{i=1}^\infty A_i \right) = \sum_{i=1}^\infty \PP(A_i)
    \]
\end{itemize}
}

\proppr{measure_properties}{
We have the following results :
    \begin{enumerate}
        \item If $A \subset B$ then : \[ \PP(B\setminus A) = \PP(B) - \PP(A) \]
        \item If $A,B \in \A$  then : \[ \PP(A \cup B) = \PP(A) + \PP(B) - \PP(A \cap B) \]
        \item If $A_n \in \A$ and $A_n \subset A_{n+1}$ for every $n\in\NN$ then : \[ \PP(\bigcup_{n\in\NN} A_n ) = \lim_{\goto{n}} \PP(A_n)\]
        \item If $A_n \in \A$ and $A_{n+1} \subset A_{n}$ for every $n\in\NN$ then : \[ \PP(\bigcap_{n\in\NN} A_n ) = \lim_{\goto{n}} \PP(A_n)\]
        \item If $A_n \in \A$  for every $n\in\NN$ then : \[ \PP(\bigcup_{n\in\NN} A_n ) \le \sum_{n\in\NN} \PP(A_n)\]
    \end{enumerate}
}{ 
    Left exercise for the reader.
}

\subsection{Monotone class}
In this section, we state and prove the monotone class theorem, which is a
fundamental tool in probability theory.

\defnr{Monotone class}{Monotone_class}{
Let $\Omega$ be a set and let $\M \subset \P(\Omega)$.   The collection $\M$ is called a \emph{monotone class} if:
\begin{enumerate}
    \item $\Omega \in \M$;
    \item if $A,B \in \M$ and $A \subset B$, then $B \setminus A \in \M$;
    \item if $(A_n)_{n \in \NN}$ is an increasing sequence in $\M$, i.e.
    \[ A_1 \subset A_2 \subset \cdots, \]
    then
    \[ \bigcup_{n \in \NN} A_n \in \M.\]
\end{enumerate}
}
Every \sfield is a monotone class. Conversely, a monotone class is a \sfield if and only if it is closed under finite intersections. Indeed, closure under finite intersections implies closure under complements and finite unions, and property (3) then yields closure under countable unions.

Arbitrary intersections of monotone classes are again monotone classes. Given a collection $\C \subset \P(\Omega)$, we define the \emph{monotone class generated by $\C$}, denoted $\M(\C)$, as the intersection of all monotone classes containing $\C$.

\thmpr{Monotone Class Theorem}{monotone_class_theorem}{
    Let $\C \subset \P(\Omega)$ be a collection of sets that is closed under finite intersections. Then
    \[\M(\C) = \sigma(\C).\]
    Consequently, if $\M$ is any monotone class such that $\C \subset \M$, then
    \[\sigma(\C) \subset \M.\]
}{
It is enough to prove the first assertion. Since any \sfield is a monotone class, it is clear that $\M(C) \subset \sigma(\C)$. To prove the reverse inclusion, it is enough to verify that $\M(\C)$ is a \sfield. As explained above, it suffices to verify that $\M(\C)$ is closed under finite intersections.

For every $A \in \P(\Omega)$, set
\[ \M_A = \{B \in \M(\C) : A \cap B \in \M(\C)\}. \]

Fix $A \in \C$. Since $\C$ is closed under finite intersections, we have $\C \subset \M_A$. Let us verify that $\M_A$ is a monotone class:

\begin{itemize}
\item $\Omega \in \M_A$ is trivial.
\item If $B, B' \in \M_A$ and $B \subset B'$, then $A \cap (B' \setminus B) = (A \cap B') \setminus (A \cap B) \in \M(\C)$ and thus $B' \setminus B \in \M_A$.
\item If $B_n \in \M_A$ for every $n \geq 0$ and the sequence $(B_n)_{n \geq 0}$ is increasing, we have $A \cap (\cup_{n \in \NN} B_n) = \cup_{n \in \NN} (A \cap B_n) \in \M(\C)$ and therefore $\cup_{n \in \NN} B_n \in \M_A$.
\end{itemize}

Since $\M_A$ is a monotone class that contains $\C$, $\M_A$ also contains $\M(\C)$. We have thus obtained that
\[
\forall A \in \C, \ \forall B \in \M(\C), \ A \cap B \in \M(\C).
\]

This is not yet the desired result, but we can use the same idea one more time. Precisely, we now fix $A \in \M(\C)$. According to the first part of the proof, $\C \subset \M_A$. By exactly the same arguments as in the first part of the proof, we get that $\M_A$ is a monotone class. It follows that $\M(\C) \subset \M_A$, which shows that $\M(\C)$ is closed under finite intersections, and completes the proof.
}


\proppr{probability_uniqueness}{
Let $(\measspace)$ be a measurable space and let $\PP$ and $\QQ$ be two probability measures on $(\measspace)$. Suppose there exists a collection $\C \subset \A$ such that:
\begin{itemize}
    \item $\C$ is closed under finite intersections;
    \item $\sigma(\C) = \A$;
    \item $\PP(A) = \QQ(A)$ for all $A \in \C$.
\end{itemize}
Then: We have $\PP = \QQ$;
}{
Let $\G = \{A \in \A : \PP(A) = \QQ(A)\}$. By assumption, $\C \subset \G$. On the other hand, it is easy to verify that $\G$ is a monotone class. For instance, if $A, B \in \G$ and $A \subset B$, we have $\PP(B \setminus A) = \PP(B) - \PP(A) = \QQ(B) - \QQ(A) = \QQ(B \setminus A)$, and hence $B \setminus A \in \Omega$ (note that we here use the fact that both $\PP$ and $\QQ$ are finite).

If follows that $\G$ contains $\M(\C)$, which is equal to $\sigma(\C)$ by the monotone class theorem. By our assumption $\sigma(\C) = \A$, we get $\G = \A$, which means that $\PP = \QQ$.
}


\subsection{Random variables}
\defnr{Random variable}{random_variable}{
A function from $\Omega$ to $\RR^n$ is called a random variable (in measure theory is called measurable function), meaning an application $\bX$ such that:
\[ \bX^{-1}(B) \in \A, \text{ for each set } B \in \B(\RR^n) \]
where $\B(\RR^n)$ denotes the Borel \footnote{The Borel \sfield on $\RR^n$ is generated by the open subsets of $\RR^n$.  The Lebesgue measure on $\RR^n$ can be first defined on rectangles of the form $[a_1,b_1]\times\cdots\times[a_n,b_n]$ by $(b_1-a_1)\cdots(b_n-a_n)$, and then uniquely extended to a measure on the Borel \sfield. For more details, see the chapter \parencite[chap.~3]{Probability}.} \sfield. }

Based on this definition, we can define a probability measure in the probabilistic space $(\RR^n, \B(\RR^n), \PP_\bX)$, where the probability measure (also called the distribution of $\bX$) $\PP_\bX$ is defined as follows:
\[ \PP_\bX(B) = \PP(\bX^{-1}(B)) = \PP(\{ \bX \in B \}), \text{ for each set } B \in \B(\RR^n) \]
We call $\PP_\bX$ the law (or the distribution) of the random variable $\bX$. In order to characterize the distribution of a random variable, we use cumulative distribution functions and probability density functions.

\defnr{Cumulative distribution function}{cumulative_distribution_function}{
Let $\bX$ be a random variable with value space $\RR^n$. For each $\bx \in \RR^n$, we define the \textbf{cumulative distribution function} as follows:
\[ F_\bX(\bx) = \PP(\bX \leq \bx) =   \PP(X_1 \leq x_1, \cdots, X_n \leq x_n)  \]
}

\defnr{Probability density function}{density_function}{
Let $\bX$ be a random variable with value space $\RR^n$. For each $\bx \in \RR^n$, we define the \textbf{probability density function} as follows:
\[ f_\bX(\bx) = \frac{\partial^n}{\partial \bx} F_\bX(\bx) = \frac{\partial^n}{\partial x_1 \cdots \partial x_n} F_\bX(\bx) \]
if it exists on a non-negligible set of $\RR^n$. This function is a density for the probability measure $\PP_\bX$ such that:
\[ \PP_\bX(B) = \int_B d\PP_\bX =  \int_B f_\bX(\bx) d\lambda(\bx) \]
Where $\lambda$ is the Lebesgue measure on $\B(\RR^n)$
}
Finally, we have the following theorem, which ensures that a constructed variable is indeed a random variable.

\thmpr{Operations on random variables}{random_variables_operations}{
\begin{enumerate}
    \item Let $X_1,X_2$ be two random variables on $\RR$, and $f : \RR^2 \mapsto \RR$ a measurable function (w.r.t Lebesgue measure and Borel \sfield), then  $(X_1,X_2), \quad f(X_1,X_2)$ are random variables.
    \item Let $(X_n)$ be a sequence of random variables that map into $\bar{\RR}$ then :
    \[ \sup _{n \in \NN} X_n, \quad \inf_{n \in \NN} X_n, \quad \limsup _{\goto{n}} X_n, \quad  \liminf_{\goto{n}} X_n \]
    are also random variables. In particular, if the sequence $(X_n)$  converges pointwise, its limit is a random variable. 
    \item Let $(X_n)$ be a sequence of random variables that map into $\bar{\RR}$ then :
    The set $A$ of all $x \in \Omega$ such that $X_n(\omega)$ converges in $\RR$ as $\goto{n}$ is measurable ($A\in\A$). Furthermore, the function $h: \Omega \longrightarrow \RR$ defined by
    \[
    Y(\omega):= \begin{cases}
        \lim _{\goto{n}} X_n(\omega) & \text { if } \omega \in A, \\
        0 & \text { if } \omega \notin A,
                \end{cases}
    \]
    is a random variable.
\end{enumerate}
}{
The proofs could be found in \parencite[sec.~1.3]{Probability}. 
}
Finally we have this two lemmas, that would be used frequently in most of probability-propositions proofs.


\lempr{Approximation by Simple Random Variables}{approximation_simple_random_variables}{
Let \( X \) be a non-negative random variable with values in $\RR$. Then there is a sequence \( (\varphi_n )\) of simple random variables with \(\varphi_{n+1} \geq \varphi_n\) such that \( X = \lim \varphi_n \) at each point of \( \Omega \).
Where a simple random variable is a random variable having the form : 
\[
    X = \sum_{i=1}^m {c}_i \1_{E_i}
\]
For some integer $m$ and real positive values $c_i$. 
}{
The proof is left exercise for reader with the hint to use : 
\[
E_{n,k} = \{ x : k2^{-n} \leq X(\omega) < (k + 1)2^{-n} \}, \text{ and the sequence } \varphi_n = 2^{-n} \sum_{k=0}^{2^{2n}} k \1_{E_{n,k}}.
\]
}


\lempr{Doob-Dynkin Lemma}{Doob_Dynkin}{
If \( \bY\) with values in $\RR^m$ is \( \sigma(\bX) \)-measurable where $\bX$ is a random variable with values in $\RR^n$, there is a measurable function \( f : \RR^n \to \RR^m \) such that \( \bY= f( \bX) \).
}{
We'll proof the theorem for real valued random variables (with values in $\RR$). 
\begin{enumerate}

    \item \textbf{ Step 1: Indicator Functions} 
    First, consider the case where \( Y = \1_A \) for some \( A \in \sigma(X) \). 
    Since \( A \in \sigma(X) \), by definition, there exists a Borel set \( B \in \B(\RR) \) such that \( A = X^{-1}(B) \).
    Define the function \( f : \RR \to \RR \) by
    \[ f(x) = \1_B(x). \]
    Then, for any \( \omega \in \Omega \),
    \[ Y(\omega) = \1_A(\omega) = \1_{\bX^{-1}(B)}(\omega) = \1_B(X(\omega)) = f(X(\omega)). \]
    Hence, \( Y = f(X) \) for this case.

    \item \textbf{ Step 2: Simple Functions}
    Now, consider \( Y \) as a simple function, i.e., \( Y = \sum_{i=1}^n c_i \1_{A_i} \) where \( c_i \in \RR \) and \( A_i \in \sigma(X) \).
    For each \( A_i \in \sigma(X) \), there exists a Borel set \( B_i \in \B(\RR) \) such that \( A_i = X^{-1}(B_i) \).
    Define \( f : \RR \to \RR \) by
    \[ f(x) = \sum_{i=1}^n c_i \1_{B_i}(x). \]
    Then, for any \( \omega \in \Omega \),
    \[ Y(\omega) = \sum_{i=1}^n c_i \1_{A_i}(\omega) = \sum_{i=1}^n c_i \1_{\bX^{-1}(B_i)}(\omega) = \sum_{i=1}^n c_i \1_{B_i}(X(\omega)) = f(X(\omega)). \]
    Thus, \( Y = f(X) \) for this case.

    \item \textbf{ Step 3: Positive Measurable Functions}
    Consider \( Y \) as a positive \( \sigma(X) \)-measurable function. There exists a sequence of simple functions \( Y_n \) such that \( Y_n \nearrow Y \). (use of the lemma \ref{lem:approximation_simple_random_variables})
    From Step 2, for each \( n \), there exists a Borel-measurable function \( f_n : \RR \to \RR \) such that
    \[ Y_n = f_n(X). \]
    Since \( Y_n \nearrow Y \), we have \( f_n(X) \to Y \) pointwise almost everywhere. 
    Define \( f(x) = \lim_{\goto{n}} f_n(x) \). Since each \( f_n \) is Borel-measurable, and the pointwise limit of Borel-measurable functions is Borel-measurable, \( f \) is Borel-measurable.
    Therefore, \( Y = f(X) \).

    \item \textbf{ Step 4: General Measurable Functions}
    Finally, consider \( Y \) as an arbitrary \( \sigma(X) \)-measurable function. We can decompose \( Y \) into its positive and negative parts:
    \[ Y = Y^+ - Y^-, \]
    where \( Y^+ = \max(Y, 0) \) and \( Y^- = -\min(Y, 0) \).
    Both \( Y^+ \) and \( Y^- \) are positive \( \sigma(X) \)-measurable functions. By Step 3, there exist Borel-measurable functions \( f^+ \) and \( f^- \) such that
    \[ Y^+ = f^+(X) \quad \text{and} \quad Y^- = f^-(X). \]
    Define \( f : \RR \to \RR \) by
    \[ f(x) = f^+(x) - f^-(x). \]
    Then,
    \[ Y = Y^+ - Y^- = f^+(X) - f^-(X) = f(X). \]
    Thus, \( Y = f(X) \) for some Borel-measurable function \( f \).
\end{enumerate} }
A large number of proofs in probability theory follow the same general strategy. The idea is to first establish a result for very simple random variables and then extend it step by step to more general cases. This approach reflects the way many objects in probability theory are constructed from simpler ones.

More precisely, proofs often proceed through the following steps:

\begin{enumerate}
\item \textbf{Step 1: Indicator Functions.}
The result is first verified for random variables of the form 
$ \1_A $ where $A\in \A$

\item \textbf{Step 2: Simple Functions.}  
The result is then extended to simple random variables, which are finite linear combinations of indicator functions.

\item \textbf{Step 3: Positive Measurable Functions.}  
Using monotone approximation by simple functions, the result is extended to non-negative measurable random variables via the approximation lemma \ref{lem:approximation_simple_random_variables}.

\item \textbf{Step 4: General Measurable Functions.}  
Finally, the result is extended to general measurable random variables by decomposing them into their positive and negative parts.
\end{enumerate}
In many situations, some of these steps may be omitted when they follow directly from previous arguments. Nevertheless, this four-step method provides a common structure that appears repeatedly in probability theory.

\section{Expectation and Moments of Random Variables}
\subsection{Expectation}
\defnr{Expectation}{Expectation}{
Let $\bX$ be a random variable with value space $\RR^n$. The \textbf{expectation} or \textbf{expected value} of a random variable $\bX$ with respect to a probability measure $\PP$ is defined as
\[
\EE[\bX] = \int_\Omega \bX \, d\PP =  \int_{\RR^n} \bx \, d\PP_\bX = \int_{\RR^n} \bx f_\bX(\bx) d\bx
\]
The integral here is a Lebesgue integral\footnote{The Lebesgue integral is the generalisation of integration on the measurable spaces, for the rigorous definition check the chapter \parencite[chap.~2]{Probability}} with respect to the probability measure. We require the integral to exist, and the last equality holds only if a density exists. }

We recall some properties of expectations as follow :
Let \( X \) and \( Y \) be random variables, and let \( a \) and \( b \) be constants. The main properties of expectations are:
\begin{enumerate}
    \item \textbf{Probability representation:}
    \begin{equation*}
    \PP[X\in A] = \EE(\1_{X\in A})
    \end{equation*}

    \item \textbf{Linearity:}
    \begin{equation*}
    \EE[a\bX+ b\bY] = a\EE[\bX] + b\EE[\bY]
    \end{equation*}

    \item \textbf{Expectation of a Constant:}
    \begin{equation*}
    \EE[a] = a
    \end{equation*}

    \item \textbf{Non-negativity:}
    If \( \bX\geq 0 \) almost surely, then
    \begin{equation*}
    \EE[\bX] \geq 0
    \end{equation*}

    \item \textbf{Finiteness:}
    \[ \EE[X] < \infty, \Longrightarrow  X < \infty \quad \PP\text{-a.s.} \]

    \item \textbf{Zero Expectation:}
    If \( X \) is a nonnegative random variable and
    \[ \EE[X] = 0 \;\Longleftrightarrow\; X = 0 \quad \PP\text{-a.s.} \]

    \item \textbf{Equality of Expectations:}
    If $Y$ is another nonnegative random variable and
    \[ X = Y \quad \PP\text{-a.s.} \]
    then
    \[ \EE[X] = \EE[Y]. \]

    \item \textbf{Law of the Unconscious Statistician:}
    If \( g \) is a measurable function and \( \bX\) is a random variable, then
    \begin{equation*}
    \EE[g(\bX)] = \int g(\bX)d\PP 
    \end{equation*}

    \item \textbf{Fatou's Lemma:}
    If \( (X_n) \) is a sequence of positive real random variables, then
    \begin{equation*}
    \EE[\liminf_{\goto{n}} X_n] \leq \liminf_{\goto{n}} \EE[X_n]
    \end{equation*}

    \item \textbf{Monotone Convergence Theorem:}
    If \( (X_n) \) is a sequence of positive real random variables that converges to \(X\) almost surely and the sequence is a.s. increasing, then
    \begin{equation*}
    \EE[\lim_{\goto{n}} X_n] = \lim_{\goto{n}} \EE[X_n]
    \end{equation*}

    \item \textbf{Dominated Convergence Theorem:}
    If \( (X_n) \) is a sequence of real random variables that converges to \( X\) almost surely and is dominated by an integrable random variable \( Y \) (i.e., \( |X_n| \leq Y \) almost surely for all \( n \)), then
    \begin{equation*}
    \EE[\lim_{\goto{n}} X_n] = \lim_{\goto{n}} \EE[X_n]
    \end{equation*}
\end{enumerate}
Before completing the section we highly invit the reader to consult the section of $\L^p$ spaces \ref{sec:Lp spaces}


\subsection{Variance}
\defnr{Variance -- Covariance -- Correlation}{second_moments}{
Let $X \in L_\RR^2(\probspace)$ be a real-valued random variable.
\begin{enumerate}
\item \textbf{Variance.}  
The \textbf{variance} of the random variable $X$ measures how much $X$ varies around its mean. It is defined by
\[ \mathrm{Var}(X) = \EE\big[(X - \EE[X])^2\big]. \]
The \textbf{standard deviation} of $X$ is defined as
\[ \sigma(X) = \sqrt{\mathrm{Var}(X)}. \]

\item \textbf{Covariance.}  
Let $X$ and $Y$ be two random variables in $L^2(\probspace)$.  
The \textbf{covariance} of $X$ and $Y$ measures how the two variables vary together and is defined by
\[ \mathrm{Cov}(X,Y) = \EE\big[(X - \EE[X])(Y - \EE[Y])\big]
= \EE[XY] - \EE[X]\EE[Y]. \]

\item \textbf{Correlation.}  
The \textbf{correlation} between $X$ and $Y$ is a normalized version of the covariance and is defined by
\[ \rho(X,Y) = \frac{\mathrm{Cov}(X,Y)}{\sigma(X)\sigma(Y)}. \]
It takes values between $-1$ and $1$.

\item \textbf{Covariance matrix.}  
Let $\bX = (X_1,\dots,X_n)$ be a random vector in $\RR^n$ with $\bX \in L_{\RR^n}^2(\probspace)$. The \textbf{covariance matrix} of $\bX$ is defined by
\[
\mathrm{Var}(\bX) 
= \EE\big[(\bX - \EE[\bX])(\bX - \EE[\bX])^T\big]
= \big(\mathrm{Cov}(X_i,X_j)\big)_{1 \le i,j \le n}.
\]
This matrix belongs to $M_n(\RR)$.
\end{enumerate}
}

\thmpr{Markov Inequalities}{Markov_Inequalities}{
\textbf{Markov's Inequality} If $X$ is a nonnegative random variable and $a>0$,
\[ \PP(X \geq a) \leq \frac{1}{a} \EE[X] \] 
\textbf{Bienaymé-Chebyshev Inequality} If $X \in L^2(\probspace)$ and $a>0$,
\[\PP(|X-\EE[X]| \geq a) \leq \frac{1}{a^2} \operatorname{var}(X) \]
                }{
    A classical proof, left exercise for the reader.\\
    Hint : Use $\EE[\1_{X \geq a}] \leq \EE[X/a]$ for Markov's inequality, and apply Markov's inequality to the nonnegative random variable $(X-\EE[X])^2$ for Chebyshev's inequality.
                }

\subsection{Characteristic Function}
We start with a basic definition. We use the notation $x \cdot y$ for the standard Euclidean scalar product in $\RR^d$.

\defnr{Characteristic function}{8.14}{
Let $X$ be a random variable with values in $\RR^d$. The characteristic function of $X$ is the function $\Phi_X : \RR^d \longrightarrow \CC$ defined by
\[ \Phi_X(\xi) := \EE[\exp(i \xi \cdot X)], \qquad \xi \in \RR^d. \]
}

The definition of expectation \ref{defn:Expectation} allows us to write
\[ \Phi_X(\xi) = \int_{\RR^d} e^{i \xi \cdot x} \,d\PP_X(x) \]

and we can also view $\Phi_X$ as the Fourier transform of the law $\PP_X$ of $X$. For this reason, we sometimes write $\Phi_X(\xi) = \widehat{\PP_X}(\xi)$. It's easy to notice (in the case $d = 1$), the dominated convergence theorem ensures that the function $\Phi_X$ is (bounded and) continuous on $\RR^d$.

Our first goal is to show that the characteristic function determines the law $\PP_X$ of $X$ (as the name suggests!). This is equivalent to showing that the Fourier transform acting on probability measures on $\RR^d$ is injective. We begin with an important calculation for a special case of the Gaussian distribution. We will study this case in more detail in Section \ref{sec:gaussian_sec}, where we will derive the following formula:
\[ \Phi_X(\xi)=\exp \left(\mathrm{i} m \xi-\frac{\sigma^2 \xi^2}{2}\right), \quad \xi \in \RR \]

\thmpr{Characteristic function determines the law}{8.16}{
The characteristic function of a random variable $X$ with values in $\RR^d$ determines the law of this random variable. In other words, the Fourier transform acting on probability measures on $\RR^d$ is injective.
}{
Let us consider the case $d = 1$. For every $\sigma > 0$, let $g_\sigma$ be the density of the Gaussian $\N(0,\sigma^2)$ distribution:
\[
g_\sigma(x) = \frac{1}{\sigma \sqrt{2\pi}} \exp\left(-\frac{x^2}{2\sigma^2}\right), \qquad x \in \RR.
\]

Let $\mu$ be a probability measure on $\RR$, and set
\[
\begin{cases}
f_\sigma(x) &= \int_{\RR} g_\sigma(x - y)\,\mu(dy) \stackrel{(\mathrm{def})}{=} g_\sigma * \mu(x), \\ 
d\mu_\sigma(x) &= f_\sigma(x)\,dx.
\end{cases}
\]

Let $C_b(\RR)$ denote the space of all bounded continuous functions $\varphi : \RR \longrightarrow \RR$, and recall that a probability measure $\nu$ on $\RR$ is characterized by the values of $\int \varphi(x)\,d\nu(x)$ for $\varphi \in C_b(\RR)$. The result of the theorem (when $d = 1$) will follow if we are able to prove that:
\begin{enumerate}
\item $\mu_\sigma$ is determined by $\widehat{\mu}$.
\item For every $\varphi \in C_b(\RR)$, $\int \varphi(x)\mu_\sigma(dx) \longrightarrow \int \varphi(x)\mu(dx)$ as $\sigma \to 0$.
\end{enumerate}

In order to establish (1), we use the formula of the characteristic function for gaussian distribution to write, for every $x \in \RR$,
\[ \sigma \sqrt{2\pi}\, g_\sigma(x) = \exp\left(-\frac{x^2}{2\sigma^2}\right) = \int_{\RR} e^{i\xi x} g_{1/\sigma}(\xi)\, d\xi. \]

It follows that
\begin{align*}
f_\sigma(x) 
    &= \int_{\RR} g_\sigma(x - y)\,\mu(dy) \\
    &= (\sigma \sqrt{2\pi})^{-1} \int_{\RR} \left( \int_{\RR} e^{i\xi(x-y)} g_{1/\sigma}(\xi)\, d\xi \right)\mu(dy) \\
    &= (\sigma \sqrt{2\pi})^{-1} \int_{\RR} e^{i\xi x} g_{1/\sigma}(\xi)
    \left( \int_{\RR} e^{-i\xi y}\mu(dy) \right)d\xi \\
    &= (\sigma \sqrt{2\pi})^{-1} \int_{\RR} e^{i\xi x} g_{1/\sigma}(\xi)\,\widehat{\mu}(-\xi)\, d\xi.
\end{align*}

In the penultimate equality, we used the Fubini-Lebesgue theorem \ref{thm:Fubini_Lebesgue}. The application of this theorem is easily justified since $\mu$ is a probability measure and the function $g_{1/\sigma}$ is integrable with respect to Lebesgue measure. The formula of the last display shows that $f_\sigma$ (and hence $\mu_\sigma$) is determined by $\widehat{\mu}$.
As for (2), we start by writing, for every $\varphi \in C_b(\RR)$,
\begin{align*}
\int \varphi(x)\mu_\sigma(dx)
    &= \int \varphi(x)\left(\int g_\sigma(y - x)\mu(dy)\right)dx \\
    &= \int \left( \int g_\sigma(y - x)\varphi(x)\,dx \right)\mu(dy) \\
    &= \int g_\sigma * \varphi(y)\,\mu(dy).
\end{align*}
where we have again applied the Fubini-Lebesgue theorem, with the same justification as above, also the last equality use the convolution that we will see in \ref{sec:measure_product_and_independence}. Then, we note that the properties
\[\begin{cases}
\int g_\sigma(x)\,dx = 1, \\
\lim_{\sigma \to 0} \int_{\{|x| > \varepsilon\}} g_\sigma(x)\,dx = 0, \quad \forall \varepsilon > 0.
\end{cases}\]

imply that, for every $y \in \RR$,
\[ \lim_{\sigma \to 0} g_\sigma * \varphi(y) = \varphi(y) \]
thanks to Proposition \ref{prop:dirac_convolution} (1) that we will see in the section measure product and independence \ref{sec:measure_product_and_independence}. By dominated convergence, using the fact that $|g_\sigma * \varphi| \leq \sup\{|\varphi(x)| : x \in \RR\}$, we get 
\[ \lim_{\sigma \to 0} \int \varphi(x)\mu_\sigma(dx) = \lim_{\sigma \to 0} \int g_\sigma * \varphi(y)\mu(dy) = \int \varphi(x)\mu(dx), \]
which completes the proof of (2) and of the case $d = 1$ of the theorem. \\ 
The proof in the general case is similar. We now use the auxiliary functions
\[
g_\sigma^{(d)}(x_1,\ldots,x_d) = \prod_{j=1}^d g_\sigma(x_j)
\]
noting that, for $\xi = (\xi_1,\ldots,\xi_d) \in \RR^d$,
\[
\int_{\RR^d} g_\sigma^{(d)}(x) e^{i\xi \cdot x}\, dx
= \prod_{j=1}^d \int_{\RR} g_\sigma(x_j)e^{i\xi_j x_j}\, dx_j
= (2\pi/\sigma^2)^{d/2} g_{1/\sigma}^{(d)}(\xi).
\]

The arguments of the case $d = 1$ then apply with obvious changes.
}

\proppr{taylor_expansion_characteristic_function}{
Let $X = (X_1,\ldots,X_d)$ be a random variable with values in $\RR^d$. Assume that $\EE[(X_j)^2] < \infty$ for every $j \in \{1,\ldots,d\}$. Then, $\Phi_X$ is twice continuously differentiable and
\[
\Phi_X(\xi) = 1 + i \sum_{j=1}^d \xi_j \EE[X_j]
- \frac{1}{2} \sum_{j=1}^d \sum_{k=1}^d \xi_j \xi_k \EE[X_j X_k]
+ o(|\xi|^2)
\]
when $\xi = (\xi_1,\ldots,\xi_d)$ tends to $0$.
}{
By differentiating under the integral sign (easy application of dominance convergence theorem), we get, for every $1 \leq j \leq d$,
\[ \frac{\partial \Phi_X}{\partial \xi_j}(\xi) = i\, \EE[X_j e^{i\xi \cdot X}]. \]

The justification is easy since $|i X_j e^{i\xi \cdot X}| = |X_j|$ and $X_j \in L^2(\Omega,\A,\PP) \subset L^1(\Omega,\A,\PP)$. Similarly, for every $j,k \in \{1,\ldots,d\}$, the fact that $X_j X_k \in L^1(\Omega,\A,\PP)$ allows us to differentiate once more and to get
\[ \frac{\partial^2 \Phi_X}{\partial \xi_j \partial \xi_k}(\xi) = - \EE[X_j X_k e^{i\xi \cdot X}]. \]

Moreover, the dominated convergence theorem ensures that the quantity $\EE[X_j X_k e^{i\xi \cdot X}]$ is a continuous function of $\xi$. We have thus proved that $\Phi_X$ is twice continuously differentiable.

Finally, the last assertion of the proposition is just Taylor’s expansion of $\Phi_X$ at order two at the origin.
}

\subsection{Laplace transform}
For a random variable with values in $\RR_+$, one often uses the Laplace transform instead of the characteristic function.

\defnr{Laplace transform}{laplace_transform}{
Let $X$ be a nonnegative random variable. The Laplace transform of (the law of) $X$ is the function $L_X$ defined on $\RR_+$ by
\[ L_X(\lambda) := \EE[e^{-\lambda X}] = \int_{\RR_+} e^{-\lambda x} \,\PP_X(dx). \]
}

The dominated convergence theorem shows that $L_X$ is continuous on $[0,\infty)$. An application of the theorem of differentiation under the integral [which is just easy result of convergence theorem] sign shows that $L_X$ is infinitely differentiable on $(0,\infty)$. The function $L_X$ is also convex (as a consequence of the fact that the functions $\lambda \mapsto e^{-\lambda x}$ are convex, for every $x \geq 0$) and thus has a (possibly infinite) right derivative at $0$, which is given by
\[ L_X'(0) = -\lim_{\lambda \downarrow 0} \EE\left[\frac{1 - e^{-\lambda X}}{\lambda}\right] = -\EE[X] \]
where the second equality follows from the monotone convergence theorem.

\thmpr{Laplace transform determines the law}{law_determined_by_laplace_transform}{
If $X$ is a nonnegative random variable, the Laplace transform of $X$ determines the law $\PP_X$ of $X$.
}{
Suppose $X$ and $X'$ are two nonnegative random variables, and that $L_X = L_{X'}$. For every $\lambda \geq 0$ and every $x \geq 0$, set $\psi_\lambda(x) = e^{-\lambda x}$. We extend the functions $\psi_\lambda$ to $[0,\infty]$ by continuity, setting $\psi_\lambda(\infty) = 0$ if $\lambda > 0$, and $\psi_0(\infty) = 1$.

Let $H$ be the vector space spanned by the functions $\psi_\lambda$, $\lambda \geq 0$. Then $H$ is a subspace of the space $C([0,\infty], \RR)$ of all real continuous functions on $[0,\infty]$, which is equipped with the topology of uniform convergence. Moreover, the Stone-Weierstrass theorem ensures that $H$ is dense in $C([0,\infty], \RR)$.

The equality $L_X = L_{X'}$ implies that $\EE[\phi(X)] = \EE[\phi(X')]$ for every $\phi \in H$. On the other hand, the mapping $\phi \mapsto \EE[\phi(X)] = \int \phi \, d\PP_X$ is clearly continuous on $C([0,\infty], \RR)$. Since $H$ is dense, it follows that we have $\EE[\psi(X)] = \EE[\psi(X')]$ for every $\psi \in C([0,\infty], \RR)$, and in particular $\EE[\psi(X)] = \EE[\psi(X')]$ for every continuous function $\psi : \RR_+ \to \RR$ with compact support, which suffices to say that $\PP_X = \PP_{X'}$.
}

Finally, in the case of integer-valued random variables, one often uses generating functions.

\defnr{Generating function}{Generating_function}{
Let $X$ be a random variable with values in $\ZZ_+$. The generating function of $X$ is the function $g_X : [0,1] \to [0,1]$ defined by
\[
g_X(r) := \EE[r^X] = \sum_{n=0}^{\infty} \PP(X = n)\, r^n.
\]
}

Note that $g_X(0) = \PP(X = 0)$ and $g_X(1) = 1$. The function $g_X$ is continuous on $[0,1]$. Since the radius of convergence of the series in the definition is at least $1$, the function $g_X$ is infinitely differentiable on $[0,1)$. Moreover, it is obvious that $g_X$ determines $\PP_X$, since the quantities $\PP(X = n)$ appear in the Taylor expansion of $g_X$ at the origin.

The derivative of $g_X$ is
\[ g_X'(r) = \sum_{n=1}^{\infty} n\, \PP(X = n)\, r^{n-1} \]
for $r \in [0,1)$. It follows that $g_X$ has a (possibly infinite) left-derivative at $r = 1$, and
\[ g_X'(1) = \sum_{n=1}^{\infty} n\, \PP(X = n) = \EE[X]. \]

More generally, for every integer $p \geq 1$, if $g_X^{(p)}$ denotes the $p$-th derivative of $g_X$,
\[\lim_{r \uparrow 1} g_X^{(p)}(r) = \EE[X(X-1)\cdots(X-p+1)]\]
which shows how to recover all moments of $X$ from the knowledge of its generating function.

\section{Lp Spaces} \label{sec:Lp spaces}
This section is mainly devoted to the study of the Lebesgue space $\L^p$ of random variables whose $p$-th power of the absolute value is having expectaiton (integrable) on a given probability measure. The fundamental inequalities of Hölder, Minkowski and Jensen provide essential ingredients for this study. We investigate the Banach space structure of $\L^p$ , and in the special case $p = 2$, the Hilbert space structure of $\L^2$, which has important applications in probability theory. 

Density theorems showing that any function of $\L^p$ can be well approximated by \textit{smoother} functions play an important role in many analytic developments. As an application of the Hilbert space structure of $L^2$ , we establish the Radon-Nikodym theorem, which allows one to decompose an arbitrary measure as the sum of a measure having a density with respect to the reference measure and a singular measure. The Radon-Nikodym theorem will be a crucial ingredient of the theory of conditional expectations.


\subsection{Definition of the space}
Throughout this section, we consider a probability space $(\probspace)$. For a real $p \geq 1$, we let $\L^p_{\RR^d}(\probspace)$ denote the space of all random variables $\bX: \Omega  \mapsto \RR^d$ such that:
\[ \int|\bX|^p \mathrm{~d} \PP<\infty  \]
We also introduce the space $\L^{\infty}_{\RR^d}(\probspace)$ of all random variables $\bX: \Omega \rightarrow \RR$ such that there exists a constant $C \in \RR_{+}$with
\[ |\bX| \leq C, \PP \text { a.e. } \]

For every $p\in[1,\infty]$, we define an equivalence relation on the set of
$p$-integrable random variables by setting : 
\[ \bX \equiv \bY \quad \text{if and only if} \quad \bX=\bY,\ \PP\text{-a.s.} \]

We then consider the quotient space
\[ L^p_{\RR^d}(\probspace)= \L^p_{\RR^d}(\probspace)/\equiv, \]

An element of $L^p_{\RR^d}(\probspace)$ is thus an equivalence class of random variables that are equal almost surely. In fact, $\conj{X}\in\L^p_{\RR^d}(\probspace)$ is defined as : 
\[ \conj{X} = \left\{ Y \in \L^p_{\RR^d}(\Omega,\A,\PP)\quad \text{such that } \bX = \bY \quad \PP-a.s. \right\}\] 
In fact we can define operation 


In what follows, we will almost always abuse notation and identify an element of $L^p_{\RR^d}(\probspace)$ with one of its representatives: when we speak of a random variable in $L^p_{\RR^d}(\probspace)$, we mean that a particular representative has been chosen, and all subsequent considerations do not depend on this choice.

We will often write $L^p(\probspace)$, $L^p(\PP)$, $L^p_{\RR^d}$ or simply $L^p$, instead of $L^p_{\RR^d}(\probspace)$ when there is no risk of confusion. Notice that the space $L^0$ is the space of all random variables, the space $L^1$ is precisely the space of integrable random variables, and that the space $L^2$ is the space of square-integrable random variables, which is of particular importance in probability theory.

For every random variable $\bX$ on $(\probspace)$ and every $p\in[1,\infty)$, we set : 
\[ \|\bX\|_p = \left(\EE\!\left[|\bX|^p\right]\right)^{1/p}, \]
with the convention $\infty^{1/p}=\infty$, and
\[ \|\bX\|_\infty = \inf\{C\in[0,\infty] : |\bX|\le C,\ \PP\text{-a.s.}\}. \]

These definitions are such that $|\bX|\le\|\bX\|_\infty$ $\PP$-a.s., and $\|\bX\|_\infty$ is the smallest number in $[0,\infty]$ with this property. We observe that, if $X$ and $Y$ are two random variables such that $X=Y$ $\PP$-a.s., then $\|\bX\|_p=\|\bY\|_p$. Hence $\|\bX\|_p$ is well defined for $\bX\in L^p(\probspace)$.

It is worth noting that when dealing with random variables taking values in $\RR^d$, the absolute value can be replaced by the Euclidean norm, indeed for a for a vector  $|\bx|$ the euclidian norm is $\sqrt{x_1^1+...+x_d^2}$.



\subsection{Principale inequalities}
\defnr{Conjugate exponents}{conjugate_exponents}{
    Let $p,q\in[1,\infty]$. We say that $p$ and $q$ are \emph{conjugate exponents} if
\[ \frac{1}{p}+\frac{1}{q}=1.\]
In particular, $p=1$ and $q=\infty$
}
\thmpr{Holder Inequality}{Holder_inequality}{
Let $p$ and $q$ be conjugate exponents. Then, if $\bX$ and $\bY$ are two random variables on $(\probspace)$,
\[ \EE\!\left[|\bX\bY|\right] \le \|\bX\|_p\,\|\bY\|_q. \]
In particular, $\bX\bY\in L^1(\probspace)$ whenever $\bX\in L^p(\probspace)$ and $\bY\in L^q(\probspace)$.
}{
If $\|\bX\|_p=0$, then $\bX=0$ $\PP$-a.s., which implies $\EE[|\bX\bY|]=0$, and the inequality is trivial. We may therefore assume that $\|\bX\|_p>0$ and $\|\bY\|_q>0$. Without loss of generality, we may also assume that $\bX\in L^p(\probspace)$ and $\bY\in L^q(\probspace)$ (otherwise $\|\bX\|_p\|\bY\|_q=\infty$ and there is nothing to prove).

The case $p=1$ and $q=\infty$ is immediate, since $|\bX\bY|\le \|\bY\|_\infty |\bX|$ $\PP$-a.s., which implies :
\[ \EE[|\bX\bY|] \le \|\bY\|_\infty\,\EE[|\bX|]  = \|Y\|_\infty\,\|\bX\|_1. \]

In what follows, we assume that $1<p<\infty$ (and thus $1<q<\infty$). Let $\alpha\in(0,1)$. Then, for every $x\in\RR_+$,
\[ x^\alpha-\alpha x \le 1-\alpha. \]
Indeed, define $\varphi_\alpha(x)=x^\alpha-\alpha x$ for $x\ge0$. For $x>0$,
\[ \varphi_\alpha'(x)=\alpha(x^{\alpha-1}-1), \]
so that $\varphi_\alpha'(x)>0$ for $x\in(0,1)$ and $\varphi_\alpha'(x)<0$ for $x\in(1,\infty)$. Hence $\varphi_\alpha$ attains its maximum at $x=1$, which yields the inequality above.

Applying this inequality to $x=u/v$, where $u\ge0$ and $v>0$, we obtain
\[ u^\alpha v^{1-\alpha} \le \alpha u + (1-\alpha)v, \]
and the inequality still holds when $v=0$. We now take $\alpha=1/p$ (so that $1-\alpha=1/q$) and set :
\[ u=\frac{|\bX(\omega)|^p}{\|\bX\|_p^p}, \qquad v=\frac{|\bY(\omega)|^q}{\|\bY\|_q^q}.\]
This yields
\[ \frac{|\bX(\omega)Y(\omega)|}{\|\bX\|_p\|\bY\|_q} \le \frac{1}{p}\frac{|\bX(\omega)|^p}{\|\bX\|_p^p} + \frac{1}{q}\frac{|Y(\omega)|^q}{\|\bY\|_q^q}, \quad \PP\text{-a.s.}
\]

Taking expectations on both sides gives
\[ \frac{1}{\|\bX\|_p\|\bY\|_q}\,\EE[|\bX\bY|] \le \frac{1}{p}+\frac{1}{q} = 1, \]
which completes the proof.
}

\thmpr{Cauchy-Schwartz Inequality}{CauchySchwartz_inequality}{
If $\bX$ and $\bY$ are two random variables on $(\probspace)$, then
\[ \EE\!\left[|\bX\bY|\right] \le \left(\EE\!\left[|\bX|^2\right]\right)^{1/2} \left(\EE\!\left[|\bY|^2\right]\right)^{1/2} = \|\bX\|_2\,\|\bY\|_2. \]
}{
    Just a special case of Holder Inequality with $p=2$
}

\proppr{Lp_spaces_inclusion}{
Suppose that $\PP$ is a probability measure and let $p>1$. Then, for any random variable $\bX\in L^p(\probspace)$,
\[ \|\bX\|_1 \le \|\bX\|_p. \]
Consequently, $L^p(\probspace)\subset L^1(\probspace)$ for every $p\in(1,\infty]$.
More generally, for any $1\le r<r'\le\infty$ and any random variable
$\bX\in L^{r'}(\probspace)$,
\[ \|\bX\|_r \le \|\bX\|_{r'}, \]
and thus $L^{r'}(\probspace)\subset L^r(\probspace)$.
}{
The inequality $\|\bX\|_1 \le \|\bX\|_p$ follows directly from Hölder's inequality by taking $\bY=1$. For the second bound in the corollary, it suffices to replace $\bX$ by $|\bX|^r$ and to take $p = r'/r$.
}   


\thmpr{Jensen Inequality}{Jensen_inequality}{
Suppose that $\PP$ is a probability measure on $(\measspace)$, and let $\varphi:\RR\to\RR_+$ be a convex function. Then, for every integrable real-valued random variable $X\in L^1_{\RR}(\probspace)$,
\[ \EE\!\left[\varphi(X)\right] \ge \varphi\!\left(\EE[X]\right). \]
\textbf{remark :} The expectation $\EE[\varphi(X)]$ is well defined, since $\varphi(X)$ is a nonnegative measurable random variable.
}{
Set
\[ \E_\varphi = \{(a,b)\in\RR^2 : \forall x\in\RR,\ \varphi(x)\ge ax+b\}. \]
Elementary properties of convex functions show that
\[ \forall x\in\RR,\quad \varphi(x)=\sup_{(a,b)\in\E_\varphi}(ax+b). \]

Since $\varphi(X)\ge aX+b$ almost surely for every $(a,b)\in\E_\varphi$, we obtain
\[ \EE[\varphi(X)] \ge \sup_{(a,b)\in\E_\varphi} \EE[aX+b] = \sup_{(a,b)\in\E_\varphi} \bigl(a\,\EE[X]+b\bigr) = \varphi\!\left(\EE[X]\right). \]   
}


\thmpr{Minkowski's Inequality}{Minkowski_inequality}{
Let $p\in[1,\infty]$, and let $\bX,\bY \in L^p(\probspace)$. Then :
\[ \bX+\bY\in L^p(\probspace) \quad \text{and} \quad \|\bX+\bY\|_p \le \|\bX\|_p + \|\bY\|_p. \]
}{
The cases $p=1$ and $p=\infty$ are immediate, using the inequality $|\bX+\bY|\le |\bX|+|\bY|$. We therefore assume that $1<p<\infty$.

Writing
\[ |\bX+\bY|^p \le (|\bX|+|\bY|)^p \le \bigl(2\max(|\bX|,|\bY|)\bigr)^p \le 2^p\bigl(|\bX|^p+|\bY|^p\bigr), \]
we see that $\EE[|\bX+\bY|^p]<\infty$, and hence $\bX+\bY\in L^p$.

Next, observe that
\[ |\bX+\bY|^p = |\bX+\bY|\cdot |\bX+\bY|^{p-1} \le |\bX|\,|\bX+\bY|^{p-1} + |\bY|\,|\bX+\bY|^{p-1}. \]
Taking expectations yields
\[ \EE[|\bX+\bY|^p] \le \EE\!\left[|\bX|\,|\bX+\bY|^{p-1}\right] + \EE\!\left[|\bY|\,|\bX+\bY|^{p-1}\right]. \]

Applying Hölder's inequality with conjugate exponents $p$ and $q=p/(p-1)$, we obtain
\[ \EE[|\bX+\bY|^p] \le \|\bX\|_p \left(\EE[|\bX+\bY|^p]\right)^{(p-1)/p} + \|\bY\|_p \left(\EE[|\bX+\bY|^p]\right)^{(p-1)/p}. \]

If $\EE[|\bX+\bY|^p]=0$, the inequality is trivial. Otherwise, dividing both sides by $\left(\EE[|\bX+\bY|^p]\right)^{(p-1)/p}$ yields
\[ \|\bX+\bY\|_p \le \|\bX\|_p + \|\bY\|_p, \]
which completes the proof.
}

\subsection{Topological structure of Lp-space}


\proppr{L0_complete}{
The formula
\[ d(X,Y) = \EE\bigl[\,|X - Y| \wedge 1\,\bigr] \]
defines a distance on $L^0_{\RR^d}(\Omega,\A,\PP)$, and the space $L^0_{\RR^d}(\Omega,\A,\PP)$ is complete for the distance $d$.
}{
Let's prove that the space $L^0_{\RR^d}(\Omega,\A,\PP)$ is complete for the distance $d$. So let $(X_n)_{n \in \NN}$ be a Cauchy sequence for the distance $d$. We can then find a subsequence $Y_k = X_{n_k}$, $k \in \NN$, such that for every $k \ge 1$,
\[d(Y_k,Y_{k+1}) \le 2^{-k}.\]

Then,
\[ \EE\Biggl[\sum_{k=1}^{\infty} \bigl(|Y_{k+1}-Y_k| \wedge 1\bigr)\Biggr] = \sum_{k=1}^{\infty} d(Y_k,Y_{k+1}) < \infty, \]
which implies that $\sum_{k=1}^{\infty} (|Y_{k+1}-Y_k| \wedge 1) < \infty$ almost surely, and therefore also $\sum_{k=1}^{\infty} |Y_{k+1}-Y_k| < \infty$ almost surely (there are almost surely only finitely many values of $k$ such that $|Y_{k+1}-Y_k| \ge 1$). We then define a random variable in $L^0_{\RR^d}(\Omega,\A,\PP)$ by setting
\[X = Y_1 + \sum_{k=1}^{\infty} (Y_{k+1}-Y_k),\]
on the event where the series converges absolutely, and $X=0$ on the complementary event, which has zero probability. By construction, the sequence $(Y_k)_{k \in \NN}$ converges almost surely to $X$, and this implies
\[ d(Y_k,X) = \EE\bigl[\,|Y_k - X| \wedge 1\,\bigr] \longrightarrow 0 \quad \text{as } k \to \infty, \]
by an application of the dominated convergence theorem. Hence the sequence $(Y_k)_{k \in \NN}$ converges to $X$ for the distance $d$, and since the sequence $(X_n)_{n \in \NN}$ is Cauchy and has a convergent subsequence, it must also converge.
}


\thmpr{Riesz Theorem}{Riesz}{ 
For every $p\in[1,\infty]$, the space $L^p(\probspace)$ equipped with the norm $\bX\mapsto \|\bX\|_p$ is a Banach space, that is, a complete normed linear space.
}{
    Can be found in  \parencite[p.~68]{Probability}
}

\thmpr{The Hilbert Space L2}{L2_Hilbert_space}{
The space $L^2_\RR(\probspace)$ equipped with the scalar product
\[ \langle X, Y \rangle = \int XY \, d\PP = \EE(XY) \]
is a real Hilbert space.
}{ The Cauchy-Schwarz inequality shows that, if $X,Y\in L^2(\probspace)$, then $XY$ is integrable and hence $\langle X,Y\rangle$ is well defined. It is then clear that the mapping $(X,Y)\mapsto \langle X,Y\rangle$ defines an inner product on $L^2(\probspace)$, and that $\langle X,X\rangle=\|X\|_2^2$. Finally, the fact that $\bigl(L^2(\probspace),\|\cdot\|_2\bigr)$ is complete is a special case of the Riesz theorem \ref{thm:Riesz}.
}
Classical results from the theory of Hilbert spaces [see appendix \ref{apen:Functional_Analysis}] can therefore be applied to $L^2(\probspace)$. In particular, if $\Phi : L^2_{\RR}(\probspace)\to\RR$ is a continuous linear functional, then there exists a unique element $Y\in L^2_{\RR^d}(\probspace)$ such that
\[ \Phi(X)=\langle X,Y\rangle \quad\text{for every } X\in L^2_{\RR^d}(\probspace), \]
where $\langle X,Y\rangle=\EE[XY]$; .

\subsection{The Radon-Nikodym Theorem}
\defnr{Measures continuity}{Measures_continuity}{
Let $\PP$ and $\QQ$ be two probability measures on $(\measspace)$. We say that:
\begin{enumerate}
\item $\QQ$ is \emph{absolutely continuous} with respect to
$\PP$ (notation $\QQ \ll \PP$) if
\[ \forall A\in\A,\quad \PP(A)=0 \;\Rightarrow\; \QQ(A)=0. \]

\item $\QQ$ is \emph{singular} with respect to $\PP$ (notation $\QQ \perp \PP$) if there exists $\Omega_0\in\A$ such that
\[ \PP(\Omega_0)=0 \quad\text{and}\quad \QQ(\Omega_0^c)=0. \]

\item $\PP$ and $\QQ$ are said to be \emph{equivalent} (notation $\PP \sim \QQ$) if
\[ \PP \ll \QQ \quad\text{and}\quad \QQ \ll \PP. \] 
\end{enumerate}
}
Now we present the theorem of Radon-Nikodym : 
\thmpr{Radon-Nikodym Theorem}{RadonNikodym}{
Suppose that $\PP$ is a $\sigma$-finite measure\footnote{A measure $\PP$ on $(\Omega,\A)$ is called \emph{finite} if $\PP(\Omega)<\infty$. It is called \emph{$\sigma$-finite} if there exists a sequence of measurable sets $(A_n)_{n\ge1}$ such that $\Omega=\bigcup_{n=1}^{\infty} A_n$ and $\PP(A_n)<\infty$ for all $n$.} on $(\measspace)$, and let $\QQ$ be another $\sigma$-finite measure on $(\measspace)$. Then there exists a unique pair $(\QQ_a,\QQ_s)$ of $\sigma$-finite measures on $(\measspace)$ such that:
\begin{enumerate}
\item $\QQ = \QQ_a + \QQ_s$,
\item $\QQ_a \ll \PP$ and $\QQ_s \perp \PP$.
\end{enumerate}
Moreover, there exists a unique\footnote{Uniqueness is in the sense that if $\tilde g$ is another measurable function satisfying the same properties, then $\tilde g = g$ $\PP$-a.s.} nonnegative measurable function $g:\Omega\to\RR_+$ such that
\[ d\QQ_a = g\cdot d\PP, \]
meaning that
\[ \forall A\in\A,\qquad \QQ_a(A)=\int_A g\,d\PP. \]
}{
    Could be found in \parencite[p.~76]{Probability} proven in the cadre of measure theory.
}
\textbf{Remark:} A useful formula one can use in \textit{change of measures} in Radon-Nikodym calculations is this : 
\[ d\QQ = g\cdot d\PP \quad\quad \textit{then} \quad\quad \int_A f d\QQ =  \int_A f\cdot g\cdot d\PP \]

\thmr{Radon-Nikodym Theorem in Probability}{RadonNikodym_prob}{
    Suppose that $\PP$ and $\QQ$ are two probability measures such that $\QQ \ll \PP$ then :\\
    There exists a unique\footnote{Uniqueness is in the sense that if $\tilde X$ is another random variable satisfying the same properties, then $\tilde g = g$ $\PP$-a.s.} nonnegative random variable $X:\Omega\to\RR_+$ such that
    \[ d\QQ = X\cdot d\PP, \]
    Furthermore, if $\PP \ll \QQ$ then $1/X$ satisfies :
    \[ d\PP = \frac{1}{X}\cdot d\QQ, \]
    
}


\proppr{equivalent_definition_absolute_continuity}{
Let $\PP$ be a $\sigma$-finite measure on $(\measspace)$, and let $\QQ$ be a finite measure on $(\measspace)$. The following conditions are equivalent:
\begin{enumerate}
\item $\QQ \ll \PP$;
\item For every $\varepsilon>0$, there exists $\delta>0$ such that, for every $A\in\A$,
\[ \PP(A)<\delta \;\Longrightarrow\; \QQ(A)<\varepsilon. \]
\end{enumerate}
}{
The implication $(\mathrm{ii}) \Rightarrow (\mathrm{i})$ is immediate.
Conversely, suppose that $\QQ \ll \PP$. By the Radon-Nikodym theorem \ref{thm:RadonNikodym}, there exists a nonnegative measurable function $g:\Omega\to\RR_+$ such that
\[ d\QQ = g\cdot d\PP. \]
Note that
\[ \int_\Omega g\,d\PP = \QQ(\Omega) < \infty. \]
The dominated convergence theorem then yields
\[ \int_\Omega \1_{\{g>n\}}\, g\, d\PP \conv{n} 0. \]
Given $\varepsilon>0$, we may therefore choose $M\in\NN$ large enough such that
\[ \int_\Omega \1_{\{g>M\}}\, g\, d\PP < \frac{\varepsilon}{2}. \]
Setting $\delta=\varepsilon/(2M)$, we obtain, for every $A\in\A$ such that $\PP(A)<\delta$,
\[ \QQ(A) = \int_A g\,d\PP \le \int_\Omega \1_{\{g>M\}}\, g\, d\PP + \int_{A\cap\{g\le M\}} g\,d\PP < \frac{\varepsilon}{2} + M\,\PP(A) \le \varepsilon. \]
This proves $(\mathrm{i}) \Rightarrow (\mathrm{ii})$ and completes the proof.
}


\section{ Measure product and Independence} \label{sec:measure_product_and_independence}
The notion of independence is one of the first genuinely new concepts introduced in probability theory, in the sense that it does not arise as a straightforward adaptation of a corresponding notion from measure theory. While the independence of two events or two random variables admits a natural intuitive interpretation, the most fundamental and structurally relevant notion is that of independence between two or more \sfields.

A central result shows that two random variables are independent if and only if the joint distribution of the pair coincides with the product of their marginal distributions. Combined with Fubini's theorem, which we present first, this characterization yields several equivalent and practical formulations of independence.

We then establish the classical Borel-Cantelli lemma and illustrate its strength through an application to the dyadic expansion of a randomly chosen real number, leading to several striking probabilistic properties.

\subsection{Probabilities measures product}
Let $(\measspace[1])$ and $(\measspace[2])$ be two probabilistic spaces. The product \sfield $\A_1\otimes\A_2$ is defined as the \sfield on $\Omega_1\times\Omega_2$ given by
\[ \A_1\otimes\A_2 = \sigma\bigl(A\times B : A\in\A_1,\ B\in\A_2\bigr). \]

Sets of the form $A\times B$, with $A\in\A_1$ and $B\in\A_2$, are called \emph{measurable rectangles}. 

\textbf{remark :} The definition of the product \sfield extends naturally to a finite collection of probabilistic spaces $(\Omega_1,\A_1),\dots,(\Omega_n,\A_n)$ by setting
\[ \A_1\otimes\cdots\otimes\A_n = \sigma(\A_1\times\cdots\times\A_n). \]
The expected associativity properties hold; for instance, if $n=3$,
\[ (\A_1\otimes\A_2)\otimes\A_3 = \A_1\otimes(\A_2\otimes\A_3) = \A_1\otimes\A_2\otimes\A_3. \]

Let $C\subset \Omega_1\times\Omega_2$ and let $\omega_1\in\Omega_1$. We define the \emph{section} of $C$ at $\omega_1$ by
\[ C_{\omega_1} = \{\omega_2\in\Omega_2 : (\omega_1,\omega_2)\in C\}. \]
Similarly, for $\omega_2\in\Omega_2$, define
\[ C^{\omega_2} = \{\omega_1\in\Omega_1 : (\omega_1,\omega_2)\in C\}. \]

If $X$ is a random variable on $\Omega_1\times\Omega_2$ and $\omega_1\in\Omega_1$, we denote by $X_{\omega_1}$ the random variable on $\Omega_2$ defined by
\[ X_{\omega_1}(\omega_2)=X(\omega_1,\omega_2). \]
Similarly, for $\omega_2\in\Omega_2$, we define $X^{\omega_2}(\omega_1)=X(\omega_1,\omega_2)$.

\proppr{product_sigma_field_sections}{
\begin{enumerate}
\item If $C\in\A_1\otimes\A_2$, then for every
$\omega_1\in\Omega_1$, the section $C_{\omega_1}$ belongs to $\A_2$,
and for every $\omega_2\in\Omega_2$, the section $C^{\omega_2}$ belongs to
$\A_1$.

\item Let $X$ be a random variable on $\Omega_1\times\Omega_2$ with respect to
$\A_1\otimes\A_2$. Then, for every $\omega_1\in\Omega_1$,
$X_{\omega_1}$ is $\A_2$-measurable, and for every
$\omega_2\in\Omega_2$, $X^{\omega_2}$ is $\A_1$-measurable.
\end{enumerate}
}{
\begin{enumerate}
\item Fix $\omega_1 \in \Omega_1$ and set
\[ \C = \{ C \in \A_1 \otimes \A_2 : C_{\omega_1} \in \A_2 \}. \]

Then $\C$ contains all measurable rectangles (if $C = A \times B$, we have $C_{\omega_1} = B$ or $C_{\omega_1} = \vide$ according as $\omega_1 \in A$ or $\omega_1 \notin A$). On the other hand, it is easy to verify that $\C$ is a $\sigma$-field, and it follows that $\C = \A_1 \otimes \A_2$. This gives the property $C_{\omega_1} \in \A_2$ for every $C \in \A_1 \otimes \A_2$. The property $C^{\omega_2} \in \A_1$ is derived in a similar manner.

\item For every measurable subset $D$ of $G$,
\begin{align*}
X_{\omega_1}^{-1}(D)
&= \{ \omega_2 \in \Omega_2 : X_{\omega_1}(\omega_2) \in D \} \\
&= \{ \omega_2 \in \Omega_2 : (\omega_1,\omega_2) \in X^{-1}(D) \} \\
&= (X^{-1}(D))_{\omega_1}.
\end{align*}

which belongs to $\A_2$ by part (1).
\end{enumerate}

}

\thmr{Existence of product measure}{product_measure_existence}{
Let $\PP_1$ and $\PP_2$ be two probability measures on $(\measspace[1])$ and $(\measspace[2])$, respectively.
\begin{enumerate}
\item There exists a unique measure $\PP_1\otimes\PP_2$ on $(\Omega_1\times\Omega_2,\A_1\otimes\A_2)$ such that
\[ (\PP_1\otimes\PP_2)(A\times B) = \PP_1(A)\,\PP_2(B), \qquad A\in\A_1,\ B\in\A_2, \]


\item For every $C\in\A_1\otimes\A_2$, the mappings
\[ \omega_1\mapsto \PP_2(C_{\omega_1}) \quad\text{and}\quad \omega_2\mapsto \PP_1(C^{\omega_2}) \]
are measurable (ie random variables), and
\[ (\PP_1\otimes\PP_2)(C) = \int_{\Omega_1}\PP_2(C_{\omega_1})\,d\PP_1(\omega_1) = \int_{\Omega_2}\PP_1(C^{\omega_2})\,d\PP_2(\omega_2). \]
\end{enumerate}
}

\textbf{Remarks:} 
\begin{enumerate}
\item Suppose now that we consider an arbitrary finite number of $\sigma$-finite measures $\PP_1,\ldots,\PP_n$, defined on $(E_1,\A_1),\ldots,(E_n,\A_n)$ respectively. We can then define the product measure
\[ \PP_1 \otimes \cdots \otimes \PP_n \]
on
\[ (E_1 \times \cdots \times E_n,\ \A_1 \otimes \cdots \otimes \A_n) \]
by setting
\[ \PP_1 \otimes \cdots \otimes \PP_n = \PP_1 \otimes (\PP_2 \otimes (\cdots \otimes \PP_n)). \]
The way we insert parentheses is in fact unimportant since $\PP_1 \otimes \cdots \otimes \PP_n$ is characterized by its values on (measurable) rectangles,
\[ \PP_1 \otimes \cdots \otimes \PP_n(A_1 \times \cdots \times A_n) = \PP_1(A_1)\cdots\PP_n(A_n). \]
\end{enumerate}


\thmpr{Fubini-Tonelli Theorem}{Fubini_Tonelli}{
Let $(\measspace[1])$ and $(\measspace[2])$ be probabilistic spaces, and equip $\Omega_1\times\Omega_2$ with the product \sfield $\A_1\otimes\A_2$ and Let $\PP_1$ and $\PP_2$ be two probability measures on $(\measspace[1])$ and $(\measspace[2])$, respectively, and let $X:\Omega_1\times\Omega_2\to[0,\infty]$ be a positive random variable.
\begin{enumerate}
\item The functions
\[\Omega_1\ni \omega_1 \longmapsto \int_{\Omega_2} X(\omega_1,\omega_2)\, d\PP_2(\omega_2), \qquad \Omega_2\ni \omega_2 \longmapsto \int_{\Omega_1} X(\omega_1,\omega_2)\, d\PP_1(\omega_1), \]
with values in $[0,\infty]$, are $\A_1$-measurable and $\A_2$-measurable, respectively (ie. they are random variables).

\item We have :
\begin{align*}
\int_{\Omega_1 \times \Omega_2} X \, d(\PP_1 \otimes \PP_2)
&= \int_{\Omega_1} \left( \int_{\Omega_2} X(\omega_1,\omega_2)\, d\PP_2(\omega_2)\right)d\PP_1(\omega_1) \\
&= \int_{\Omega_2}\left(\int_{\Omega_1}X(\omega_1,\omega_2)\, d\PP_1(\omega_1)\right)d\PP_2(\omega_2).
\end{align*}

Or more simply :
\[ \EE_{\PP_1 \otimes \PP_2}[X] = \EE_{\PP_1}\Big[
      \EE_{\PP_2}\big[ X(\omega_1,\cdot) \big] \Big] = \EE_{\PP_2}
   \Big[ \EE_{\PP_1}\big[ X(\cdot,\omega_2) \big]\Big].
\]

\end{enumerate}
}{
We give the proof for the special case $X=\1_C$ for some $C\in\A_1\otimes\A_2$. The general case can be deduced from this special case by applying the monotone convergence theorem to the sequence of random variables $X_n \conv{n} X$. \\
In fact, let $X=\1_C$ for some $C\in\A_1\otimes\A_2$ the functions in the theorem are reduced to : 
\[\Omega_1\ni \omega_1 \longmapsto \int_{\Omega_2} \1_{C_{\omega_1}}(\omega_2), d\PP_2(\omega_2), \qquad \Omega_2\ni \omega_2 \longmapsto \int_{\Omega_1} \1_{C^{\omega_2}}(\omega_1)\, d\PP_1(\omega_1), \]
Which are exactly : 
\[ \omega_1\mapsto \PP_2(C_{\omega_1}) \quad\text{and}\quad \omega_2\mapsto \PP_1(C^{\omega_2}) \]
By the second point of Theorem \ref{thm:product_measure_existence}, these functions are measurable, and we have
\[ (\PP_1\otimes\PP_2)(C) = \int_{\Omega_1}\PP_2(C_{\omega_1})\,d\PP_1(\omega_1) = \int_{\Omega_2}\PP_1(C^{\omega_2})\,d\PP_2(\omega_2). \]
Which is :
\[ \EE_{\PP_1 \otimes \PP_2}[X] = \EE_{\PP_1}\Big[
      \EE_{\PP_2}\big[ X(\omega_1,\cdot) \big] \Big] = \EE_{\PP_2}
   \Big[ \EE_{\PP_1}\big[ X(\cdot,\omega_2) \big]\Big].
\]
}

\thmpr{Fubini-Lebesgue Theorem}{Fubini_Lebesgue}{
Let $X\in L^1(\Omega_1\times\Omega_2,\A_1\otimes\A_2, \PP_1\otimes\PP_2)$. Then:
\begin{enumerate}
\item For $\PP_1$-almost every $\omega_1\in\Omega_1$, the random variable $\omega_2\mapsto X(\omega_1,\omega_2)$ belongs to $L^1(\probspace[2])$, and for $\PP_2$-almost every $\omega_2\in\Omega_2$, the random variable $\omega_1\mapsto X(\omega_1,\omega_2)$ belongs to $L^1(\probspace[1])$.

\item The functions
\[ \omega_1 \longmapsto \int_{\Omega_2} X(\omega_1,\omega_2)\, d\PP_2(\omega_2), \qquad \omega_2 \longmapsto \int_{\Omega_1} X(\omega_1,\omega_2)\, d\PP_1(\omega_1), \]
belong to $L^1(\probspace[1])$ and $L^1(\probspace[2])$, respectively.

\item We have :
\begin{align*}
\int_{\Omega_1 \times \Omega_2} X \, d(\PP_1 \otimes \PP_2)
&= \int_{\Omega_1} \left( \int_{\Omega_2} X(\omega_1,\omega_2)\, d\PP_2(\omega_2)\right)d\PP_1(\omega_1) \\
&= \int_{\Omega_2}\left(\int_{\Omega_1}X(\omega_1,\omega_2)\, d\PP_1(\omega_1)\right)d\PP_2(\omega_2).
\end{align*}

Or more simply :
\[ \EE_{\PP_1 \otimes \PP_2}[X] = \EE_{\PP_1}\Big[
      \EE_{\PP_2}\big[ X(\omega_1,\cdot) \big] \Big] = \EE_{\PP_2}
   \Big[ \EE_{\PP_1}\big[ X(\cdot,\omega_2) \big]\Big].
\]

\end{enumerate}
}{
\begin{enumerate}
\item  By theorem \ref{thm:Fubini_Tonelli} applied to $|X|$, we have
\[ \int_{\Omega_1} \left( \int_{\Omega_2} |X(\omega_1,\omega_2)| \, d\PP_2(\omega_2) \right)d\PP_1(\omega_1) = \int |X| \, d\PP_1 \otimes \PP_2 < \infty. \]

Thanks to propertry (5) following definition of expectation \ref{defn:Expectation}, this implies that
\[ \int_{\Omega_2} |X(\omega_1,\omega_2)| \, d\PP_2(\omega_2) < \infty \quad \text{for } \PP_1\text{-a.e. } \omega_1. \]
Thus, the function $\omega_1 \mapsto X(\omega_1,\omega_2)$ (which is measurable by proposition \ref{prop:product_sigma_field_sections}) belongs to $L^1(\Omega_2,\A_2,\PP_2)$ except on the $\A_1$-measurable set
\[ N := \left\{ \omega_1 \in \Omega_1 : \int_{\Omega_2} |X(\omega_1,\omega_2)| \, d\PP_2(\omega_2) = \infty \right\}, \]
which has zero $\PP_1$-measure. The same argument applies to the function $\omega_2 \mapsto X(\omega_1,\omega_2)$.

\medskip

\item  If $\omega_1 \in N^c$, $\int_{\Omega_2} X(\omega_1,\omega_2)\,d\PP_2(\omega_2)$ is well defined since the function is only defined (thanks to (1)) outside a measurable set of values of $\omega_1$ of zero $\PP_1$-measure. We can always agree that that $\int_{\Omega_2} X(\omega_1,\omega_2)\,d\PP_2(\omega_2)=0$ if $\omega_1 \in N$. Then the function
\begin{align*}
\omega_1 \mapsto \int_{\Omega_2} X(\omega_1,\omega_2)\,d\PP_2(\omega_2)
&= \1_{N^c}(\omega_1)\int_{\Omega_2} X^+(\omega_1,\omega_2)\,d\PP_2(\omega_2) \\
&\quad - \1_{N^c}(\omega_1)\int_{\Omega_2} X^-(\omega_1,\omega_2)\,d\PP_2(\omega_2).
\end{align*}

is $\A_1$-measurable by \ref{thm:Fubini_Tonelli} (1), and we have also, thanks to part (2) of the same theorem,
\begin{align*}
\int_{\Omega_1} \left| \int_{\Omega_2} X(\omega_1,\omega_2)\,d\PP_2(\omega_2) \right| d\PP_1(\omega_1)
&\leq \int_{\Omega_1} \left( \int_{\Omega_2} |X(\omega_1,\omega_2)|\,d\PP_2(\omega_2) \right)d\PP_1(\omega_1) \\
&= \int |X| \, d\PP_1 \otimes \PP_2 \\
&< \infty.
\end{align*}

This proves that the function $\omega_1 \mapsto \int_{\Omega_2} X(\omega_1,\omega_2)\,d\PP_2(\omega_2)$ is in $L^1(\Omega_1,\A_1,\PP_1)$.
The same argument applies to the function $\omega_2 \mapsto \int_{\Omega_1} X(\omega_1,\omega_2)\,d\PP_1(\omega_1)$.

\medskip

\item Theorem \ref{thm:Fubini_Tonelli} gives the two equalities
\begin{align*}
\int_{\Omega_1} \left( \int_{\Omega_2} X^+(\omega_1,\omega_2)\,d\PP_2(\omega_2) \right)d\PP_1(\omega_1)
&= \int_{\Omega_1 \times \Omega_2} X^+ \, d\PP_1 \otimes \PP_2, \\
\int_{\Omega_1} \left( \int_{\Omega_2} X^-(\omega_1,\omega_2)\,d\PP_2(\omega_2) \right)d\PP_1(\omega_1)
&= \int_{\Omega_1 \times \Omega_2} X^- \, d\PP_1 \otimes \PP_2.
\end{align*}

We can replace the integral over $\Omega_1$ by an integral over $N^c$ in the left-hand sides
of these formulas. Then we just have to subtract the second equality from the first one
(using part (2)), in order to arrive at the first equality in (3), and the second equality
is derived in a similar manner.
\end{enumerate}
}

\subsection{Convolution}

Throughout this section, we consider the measure space $(\RR^d, \B(\RR^d), \lambda)$. Let $f$ and $g$ be two real measurable functions on $\RR^d$. For $x \in \RR^d$, the convolution
\[ (f * g)(x) = \int_{\RR^d} f(x-y)g(y)\,dy \]
is well defined provided
\[ \int_{\RR^d} |f(x-y)g(y)|\,dy < \infty. \]

In that case, the fact that Lebesgue measure on $\RR^d$ is invariant under translations and under the symmetry $y \mapsto -y$ shows that $(g * f)(x)$ is also well defined and $(g * f)(x) = (f * g)(x)$.

\proppr{conv_Lp_inequ}{
Let $f \in L^1(\RR^d,\B(\RR^d),\lambda)$ and $g \in L^p(\RR^d,\B(\RR^d),\lambda)$, where $p \in [1,\infty)$. \\ 
Then, for $\lambda$ a.e. $x \in \RR^d$\footnote{ for all $x$ except a set of null measure}, the convolution $(f * g)(x)$ is well defined.\\
Moreover $f * g \in L^p(\RR^d,\B(\RR^d),\lambda)$ and : 
    \[\|f * g\|_p \leq \|f\|_1 \|g\|_p\]
}{
Again, we may agree that, on the set of zero Lebesgue measure where $f * g$ is not well defined, we take $f * g(x) = 0$.

We leave aside the case $p = \infty$, which is easier (and where stronger results are proved in Proposition \ref{prop:conv_Lpq_inequ} below). We assume that $\|f\|_1 > 0$, since otherwise the results are trivial. Since the measure with density $\|f\|_1^{-1}|f(x)|$ with respect to $\lambda$ is a probability measure, Jensen’s inequality gives, for every $x \in \RR^d$,
\begin{align*}
\left(\int_{\RR^d} |g(x-y)||f(y)|\,dy \right)^p
    &= \|f\|_1^p \left(\int_{\RR^d} |g(x-y)| \frac{|f(y)|}{\|f\|_1} dy \right)^p \\
    &\leq \|f\|_1^{p-1} \int_{\RR^d} |g(x-y)|^p |f(y)|\,dy.
\end{align*}

Using Fubini-Tonelli Theorem \ref{thm:Fubini_Tonelli}, we have then :
\begin{align*}
\int_{\RR^d} \left(\int_{\RR^d} |g(x-y)||f(y)|\,dy \right)^p dx
    &\leq \|f\|_1^{p-1} \int_{\RR^d} \left(\int_{\RR^d} |g(x-y)|^p |f(y)|\,dy \right) dx \\
    &= \|f\|_1^{p-1} \int_{\RR^d} \left(\int_{\RR^d} |g(y)|^p |f(x-y)|\,dy \right) dx \\
    &= \|f\|_1^{p-1} \int_{\RR^d} |g(y)|^p \left(\int_{\RR^d} |f(x-y)| dx \right) dy \\
    &= \|f\|_1^p \|g\|_p^p < \infty,
\end{align*}
which shows that
\[
\int_{\RR^d} |f(x-t)||g(t)|dt < \infty, \qquad \lambda(dx)\ \text{a.e.} \]


and gives the first assertion. The second one also follows since the previous calculation gives
\[
\int_{\RR^d} |f * g(x)|^p dx
\leq \int_{\RR^d} \left(\int_{\RR^d} |g(x-y)||f(y)|\,dy \right)^p dx
\leq \|f\|_1^p \|g\|_p^p.
\]
}

The next proposition gives another important case where the convolution $f * g$ is well defined and enjoys nice continuity properties.

\proppr{conv_Lpq_inequ}{
Let $p \in [1,\infty)$ and $q \in (1,\infty)$ be conjugate exponents $\left(\frac{1}{p} + \frac{1}{q} = 1\right)$. Let $f \in L^p(\RR^d,\B(\RR^d),\lambda)$ and $g \in L^q(\RR^d,\B(\RR^d),\lambda)$. Then, the convolution $f * g(x)$ is well defined and $|f * g(x)| \leq \|f\|_p \|g\|_q$, for every $x \in \RR^d$. Furthermore, the function $f * g$ is uniformly continuous on $\RR^d$.
}{
By the Hölder inequality,
\[
\int_{\RR^d} |f(x-y)g(y)|dy
\leq \left(\int |f(x-y)|^p dy \right)^{1/p} \|g\|_q
= \|f\|_p \|g\|_q.
\]

This gives the first assertion and also proves that $f * g$ is bounded above by $\|f\|_p \|g\|_q$. To get the property of uniform continuity, we use the Lemma \ref{lem:continuity_translation} it is easy to complete the proof of the proposition. Indeed, for every $x,x' \in \RR^d$,
\begin{align*}
|f * g(x) - f * g(x')|
&\leq \int |f(x-y) - f(x'-y)||g(y)|dy \\
&\leq \|g\|_q \left(\int |f(x-y) - f(x'-y)|^p dy \right)^{1/p} \\
&= \|g\|_q \|f \circ \sigma_{-x} - f \circ \sigma_{-x'}\|_p
\end{align*}
and we use Lemma \ref{lem:continuity_translation} to say that $\|f \circ \sigma_{-x} - f \circ \sigma_{-x'}\|_p$ tends to $0$ when $x-x' \to 0$.
}

\lempr{Continuity of translation}{continuity_translation}{
Set $\sigma_x(y) = y - x$ for every $x,y \in \RR^d$. Let $p \in [1,\infty)$ and $f \in L^p(\RR^d,\B(\RR^d),\lambda)$. Then the mapping $x \mapsto f \circ \sigma_x$ is uniformly continuous from $\RR^d$ into $L^p(\RR^d,\B(\RR^d),\lambda)$.
}{
The proof of this lemma is omitted, but could be found in the book \parencite[lem~5.7]{Probability}
%  Suppose first that $f$ belongs to the space $C_c(\RR^d)$ of all real continuous functions with compact support on $\RR^d$. Then,
% \[
% \|f \circ \sigma_x - f \circ \sigma_y\|_p^p
% = \int |f(z-x) - f(z-y)|^p dz
% = \int |f(z) - f(z-(y-x))|^p dz
% \]
% which tends to $0$ when $y-x \to 0$ by an application of the dominated convergence theorem. In the general case, Theorem 4.8 (3) allows us to find a sequence $(f_n)_{n \in \NN}$ in $C_c(\RR^d)$ that converges to $f$ in $L^p(\RR^d,\B(\RR^d),\lambda)$. Then, for every $x,y \in \RR^d$,
% \begin{align*}
% \|f \circ \sigma_x - f \circ \sigma_y\|_p
% &\leq \|f \circ \sigma_x - f_n \circ \sigma_x\|_p
% + \|f_n \circ \sigma_x - f_n \circ \sigma_y\|_p
% + \|f_n \circ \sigma_y - f \circ \sigma_y\|_p \\
% &= 2\|f - f_n\|_p + \|f_n \circ \sigma_x - f_n \circ \sigma_y\|_p.
% \end{align*}

% Fix $\varepsilon > 0$ and first choose $n$ large enough so that $\|f - f_n\|_p < \varepsilon/4$. Since $f_n \in C_c(\RR^d)$, we can then choose $\delta > 0$ such that $\|f_n \circ \sigma_x - f_n \circ \sigma_y\|_p \leq \varepsilon/2$ if $|x-y| < \delta$. The preceding display then shows that $\|f \circ \sigma_x - f \circ \sigma_y\|_p \leq \varepsilon$ if $|x-y| < \delta$. This completes the proof.
}


\paragraph{Approximations of the Dirac Measure}
\defnr{Approximation of the Dirac Measure}{dirac_aprox}{
Let us say that a sequence $(\varphi_n)_{n \in \NN}$ of nonnegative continuous functions on $\RR^d$ is an approximation of the Dirac measure $\delta_0$ if:
\begin{enumerate}
\item For every $n \in \NN$,
\[ \int_{\RR^d} \varphi_n(x)\,dx = 1. \]

\item For every $\delta > 0$,
\[ \lim_{n \to \infty} \int_{\{|x|>\delta\}} \varphi_n(x)\,dx = 0. \]
\end{enumerate}
}

It is easy to construct approximations of $\delta_0$. If $\varphi$ is any nonnegative continuous function on $\RR^d$ such that $\int \varphi(x)\,dx = 1$, we can set $\varphi_n(x) = n^d \varphi(nx)$ for every $x \in \RR^d$ and every $n \in \NN$. Note that, if we start from a function $\varphi$ with compact support, all functions $\varphi_n$ will vanish outside the same closed ball.

\proppr{5.8}{
Let $(\varphi_n)_{n \in \NN}$ be an approximation of the Dirac measure $\delta_0$.
\begin{enumerate}
\item For any bounded continuous function $f : \RR^d \longrightarrow \RR$, the sequence $\varphi_n * f$ converges to $f$ as $n \to \infty$, uniformly on every compact subset of $\RR^d$.
\item Let $f \in L^p(\RR^d,\B(\RR^d),\lambda)$, where $p \in [1,\infty)$. Then $\varphi_n * f \to f$ in $L^p$.
\end{enumerate}
}{
\begin{enumerate}
\item First note that $\varphi_n * f$ is well defined since $\varphi_n$ is integrable and $f$ is bounded. Then, we have, for every $\delta > 0$,
\begin{align*}
\varphi_n * f(x) - f(x)
&= \int_{|y|\leq \delta} (f(x-y) - f(x))\varphi_n(y)\,dy \\
&\quad + \int_{|y|>\delta} (f(x-y) - f(x))\varphi_n(y)\,dy. \tag{i}
\end{align*}

Fix $\varepsilon > 0$ and a compact subset $H$ of $\RR^d$. Since $f$ is uniformly continuous on any compact subset of $\RR^d$, we can choose $\delta > 0$ such that $|f(x-y)-f(x)| \leq \varepsilon$ for every $x \in H$ and $y \in \RR^d$ such that $|y| \leq \delta$. Recalling (1) in Definition \ref{defn:dirac_aprox}, we get that the first term in the right-hand side of (i) is smaller than $\varepsilon$ in absolute value, for every $x \in H$ and $n \in \NN$. Furthermore, (2) in Definition \ref{defn:dirac_aprox} and the boundedness of $f$ show that the second term in the right-hand side of (i) tends to $0$ when $n \to \infty$, uniformly when $x$ varies in $\RR^d$. It follows that we have $|\varphi_n * f(x) - f(x)| < 2\varepsilon$ for every $x \in H$, as soon as $n$ is sufficiently large.

\item The fact that $\varphi_n * f$ is well defined a.e. and belongs to $L^p$ follows from Proposition \ref{prop:conv_Lp_inequ} since $\varphi_n$ is integrable. We then observe that
\begin{align*}
\int |\varphi_n * f(x) - f(x)|^p dx
    &\leq \int \left( \int \varphi_n(y)|f(x-y)-f(x)|\,dy \right)^p dx \\
    &\leq \int \left( \int \varphi_n(y)|f(x-y)-f(x)|^p dy \right) dx \\
    &= \int \left( \int \varphi_n(y)|f(x-y)-f(x)|^p dx \right) dy \\
    &= \int \varphi_n(y)\left( \int |f(x-y)-f(x)|^p dx \right) dy
\end{align*}

where the second inequality is a consequence of Jensen's inequality (note that $\varphi_n(y)dy$ is a probability measure), and the next equality uses Fubini-Tonelli Theorem \ref{thm:Fubini_Tonelli}. For every $y \in \RR^d$, set
\[ h(y) = \int |f(x-y)-f(x)|^p dx = \|f \circ \sigma_y - f\|_p^p, \]
with the notation of Lemma \ref{lem:continuity_translation}. The function $h$ is bounded and $h(y) \to 0$ as $y \to 0$ by Lemma \ref{lem:continuity_translation}. It then immediately follows from (1) and (2) in Definition \ref{defn:dirac_aprox} that
\[ \lim_{n \to \infty} \int \varphi_n(y)h(y)\,dy = 0. \]
This completes the proof.
\end{enumerate}

}

\paragraph{Application}
In dimension $d = 1$, we can take
\[
\varphi_n(x) = c_n(1 - x^2)^n \1_{\{|x|\leq 1\}}
\]
where the constant $c_n$ is chosen so that $\int \varphi_n(x)dx = 1$. Then let $[a,b]$ be a compact subinterval of $(0,1)$, and let $f$ be a continuous function on $[a,b]$. It is easy to extend $f$ to a continuous function on $\RR$ that vanishes outside $[0,1]$ (for instance, take $f$ to be affine on both intervals $[0,a]$ and $[b,1]$). Then the preceding proposition shows that
\[
\varphi_n * f(x)
= c_n \int (1-(x-y)^2)^n \1_{\{|x-y|\leq 1\}} f(y)\,dy
\longrightarrow f(x)
\]
uniformly on $[a,b]$. If $x \in [a,b]$, we can remove the indicator function $\1_{\{|x-y|\leq 1\}}$ in the last display, since the conditions $x \in [a,b]$ and $|x-y|>1$ force $f(y)=0$. It now follows that $f$ is a uniform limit of polynomials on the interval $[a,b]$ (this is a special case of the Stone-Weierstrass theorem).



\subsection{Independent Events}

\defnr{Independence of events}{Events_independence}{
\begin{enumerate}
    \item Two events $A$ and $B$ are said to be \emph{independent} if
    \[ \PP(A\cap B)=\PP(A)\,\PP(B). \]

    \item  Let $A_1,\dots,A_n$ be events. We say that the events $A_1,\dots,A_n$ are \emph{independent} if, for every nonempty subset $\{i_1,\dots,i_p\}\subset\{1,\dots,n\}$, we have
    \[ \PP(A_{i_1}\cap\cdots\cap A_{i_p}) = \PP(A_{i_1})\cdots\PP(A_{i_p}). \]
\end{enumerate}
}
\textbf{remark:} In particular, for $n\ge3$, pairwise independence does not imply independence.

\proppr{Multi_event_independence}{
The events $A_1,\dots,A_n$ are independent if and only if
\[ \PP(B_1\cap\cdots\cap B_n) = \PP(B_1)\cdots\PP(B_n) \]
for every choice of events $B_i\in\sigma(A_i) =\{\varnothing,A_i,A_i^c,\Omega\}$, $i=1,\dots,n$.
}{
    Left exercise for the reader.
}


\subsection{Independence for \sfields and Random Variables}
We say that $\G$ is a \emph{sub-\sfield} of $\A$ if
$\G$ is a \sfield and $\G\subset\A$.
Roughly speaking, a sub-\sfield represents partial information in the
probability space, namely the information obtained by knowing which events in
$\G$ have occurred. For instance, if $X$ is a random variable, the
\sfield $\sigma(X)$ corresponds to the information provided by the value
of $X$.

This suggests that the most general notion of independence is that of
independent sub-\sfields: informally, two sub-\sfields
$\G$ and $\G'$ are independent if knowing which events of
$\G$ have occurred gives no information about the occurrence of events
in $\G'$, and conversely. We make this precise in the following
definition.

\defnr{Independence of [finite] \sfields}{sigmafield_independence}{
\begin{enumerate}
\item Let $\G_1,\dots,\G_n$ be sub-\sfields of $\A$. We say that $\G_1,\dots,\G_n$ are \emph{independent} if, for every choice of events
\[ A_1\in\G_1,\ A_2\in\G_2,\ \dots,\ A_n\in\G_n, \]
we have
\[ \PP(A_1\cap A_2\cap\cdots\cap A_n) = \PP(A_1)\,\PP(A_2)\cdots\PP(A_n). \]
\item Let $X_1,\dots,X_n$ be random variables taking values in the measurable spaces $(E_1,\E_1),\dots,(E_n,\E_n)$, respectively. We say that the random variables $X_1,\dots,X_n$ are \emph{independent} if the \sfields $\sigma(X_1),\dots,\sigma(X_n)$ are independent.

\item Equivalently, $X_1,\dots,X_n$ are independent if and only if, for every choice of sets $F_1\in\E_1,\dots,F_n\in\E_n$,
\[ \PP\bigl(X_1\in F_1,\dots,X_n\in F_n\bigr) = \PP(X_1\in F_1)\cdots\PP(X_n\in F_n). \]
\end{enumerate}
}
Firstly, we introduce a suffisant condition for the independence of \sfields. 
\thmpr{A suffisant condition for independence of \sfields}{suffisant_condition_sigmafield_independence}{
Let $\G_1,\dots,\G_n$ be sub-\sfields of $\A$. For each $i\in\{1,\dots,n\}$, let $\C_i\subset\G_i$ be a class of sets such that:
\begin{itemize}
\item $\C_i$ is closed under finite intersections,
\item $\Omega\in\C_i$,
\item $\sigma(\C_i)=\G_i$.
\end{itemize}
Assume that
\[ \forall C_1\in\C_1,\dots,C_n\in\C_n, \qquad \PP(C_1\cap\cdots\cap C_n) = \prod_{i=1}^n \PP(C_i). \]
Then the \sfields $\G_1,\dots,\G_n$ are independent.
}{
We first fix $C_2\in\C_2,\dots,C_n\in\C_n$ and define
\[ \M_1 = \Bigl\{ B_1\in\G_1 : \PP(B_1\cap C_2\cap\cdots\cap C_n) = \PP(B_1)\prod_{i=2}^n \PP(C_i) \Bigr\}. \]
By assumption, $\C_1\subset\M_1$. It is straightforward to check that $\M_1$ is a monotone class. By the monotone class theorem,
\[ \M_1 \supset \sigma(\C_1)=\G_1. \]
Hence, for all $B_1\in\G_1$ and all $C_2\in\C_2,\dots,C_n\in\C_n$,
\[ \PP(B_1\cap C_2\cap\cdots\cap C_n) = \PP(B_1)\prod_{i=2}^n \PP(C_i). \]

We now fix $B_1\in\G_1$, $C_3\in\C_3,\dots,C_n\in\C_n$ and define
\[ \M_2 = \Bigl\{ B_2\in\G_2 : \PP(B_1\cap B_2\cap C_3\cap\cdots\cap C_n) = \PP(B_1)\PP(B_2)\prod_{i=3}^n \PP(C_i) \Bigr\}. \]
Again, $\C_2\subset\M_2$, and $\M_2$ is a monotone class, so $\M_2=\G_2$.

Proceeding inductively, we obtain that, for all $B_1\in\G_1,\dots,B_n\in\G_n$,
\[ \PP(B_1\cap\cdots\cap B_n) = \prod_{i=1}^n \PP(B_i). \]
This is precisely the definition of independence of the \sfields $\G_1,\dots,\G_n$.
}

The next theorem provides a fundamental characterization of the independence of $n$ random variables.

\thmpr{Characterization of Random Variable Independence}{random_variable_independence_characterization}{
Let $X_1,\dots,X_n$ be random variables taking values in $\RR$, the following statements are
equivalent:
\begin{enumerate}
\item
The random variables $X_1,\dots,X_n$ are independent.

\item  the law of the random vector $(X_1,\dots,X_n)$ is the product of the respective laws of $X_1,\dots,X_n$, that is,
\[ \PP_{(X_1,\dots,X_n)} = \PP_{X_1}\otimes\cdots\otimes\PP_{X_n}.\]

\item For every $a_1,\dots,a_n\in\RR$,
\[ \PP(X_1\le a_1,\dots,X_n\le a_n) = \prod_{i=1}^n \PP(X_i\le a_i). \]

\item
For every choice of continuous functions $f_1,\dots,f_n:\RR\to\RR_+$ with compact support,
\[ \EE\!\left[\prod_{i=1}^n f_i(X_i)\right] = \prod_{i=1}^n \EE[f_i(X_i)]. \]

\item The characteristic function of the random vector $X=(X_1,\dots,X_n)$ factorizes as
\[ \Phi_X(\xi_1,\dots,\xi_n) = \prod_{i=1}^n \Phi_{X_i}(\xi_i), \qquad (\xi_1,\dots,\xi_n)\in\RR^n. \]
\end{enumerate}
}{
\begin{itemize}
    \item $(1)\Longleftrightarrow(2)$ : \\ 
Let $F_i \in \E_i$, for every $i \in \{1,\ldots,n\}$. On one hand,
\[ \PP_{(X_1,\ldots,X_n)}(F_1 \times \cdots \times F_n) = \PP\big( (X_1 \in F_1) \cap \cdots \cap (X_n \in F_n) \big). \]
On the other hand,
\[ \PP_{X_1} \otimes \cdots \otimes \PP_{X_n}(F_1 \times \cdots \times F_n) = \prod_{i=1}^n \PP_{X_i}(F_i) = \prod_{i=1}^n \PP(X_i \in F_i). \]
Comparing with the definition of independence, we see that $X_1,\ldots,X_n$ are independent if and only if the two probability measures $\PP_{(X_1,\ldots,X_n)}$ and $\PP_{X_1} \otimes \cdots \otimes \PP_{X_n}$ take the same values on boxes $F_1 \times \cdots \times F_n$.
This is equivalent to saying that : 
\[ \PP_{(X_1,\ldots,X_n)} = \PP_{X_1} \otimes \cdots \otimes \PP_{X_n}, \]
since we know that a probability measure on a product space is characterized by its values on boxes (theorem \ref{thm:product_measure_existence}) we conclude the result.
    \item $(2)\Longleftrightarrow(3)$ : \\ 
Just by taking $F_i=]-\infty,a_i]$ we can prove the $(2)\Longleftarrow (3)$, conversely,
Just by taking $C_i = \bigl\{ ]-\infty, a] \;\big|\; a \in \RR \bigr\}$, and apply the theorem \ref{thm:suffisant_condition_sigmafield_independence} we obtain the independence $(1)$ and therefore $(2)$.
    \item $(1)\Longleftrightarrow(4)$ : \\ 
We have $(4) \Longrightarrow (3) \Longrightarrow (1)$ by taking $f_i(x) = \1_{x\leq a_i}$.
The other implication is imediate from the next proposition. 

    \item $(2)\Longleftrightarrow(5)$ : \\  To show this equivalence, we note that the Fourier transform of the product measure $\PP_{X_1} \otimes \cdots \otimes \PP_{X_n}$ is
\begin{align*}
(\xi_1,\ldots,\xi_n)
&\longmapsto
\int \exp\!\left( i \sum_{j=1}^n \xi_j x_j \right)
\, \PP_{X_1}(dx_1) \cdots \PP_{X_n}(dx_n) \\
&= \prod_{j=1}^n \int e^{i \xi_j x_j}\, \PP_{X_j}(dx_j) \\
&= \prod_{j=1}^n \Phi_{X_j}(\xi_j).
\end{align*}

using the Fubini theorem in the first equality. Hence the property (5) is equivalent to the fact that the Fourier transform of $\PP_{(X_1,\ldots,X_n)}$ coincides with the Fourier transform of $\PP_{X_1} \otimes \cdots \otimes \PP_{X_n}$. By pervious property, this holds if and only if
\[ \PP_{(X_1,\ldots,X_n)} = \PP_{X_1} \otimes \cdots \otimes \PP_{X_n}, \]

\end{itemize}
}

From this theorem follow this interesting results :
\proppr{Random_variables_Independence_proprieties}{
Let $X_1,\dots,X_n$ be independent real-valued random variables.
\begin{enumerate}
\item
Let $f_1,\dots,f_n$ be non negative measurable functions we have : 
\[ \EE\!\left[\prod_{i=1}^n f_i(X_i)\right] = \prod_{i=1}^n \EE\!\left[f_i(X_i)\right], \]
\item
Let $f_1,\dots,f_n$ be measurable functions such that
\[ \EE[|f_i(X_i)|]<\infty \quad\text{for every } i\in\{1,\dots,n\}. \]
Then
\[ \EE\!\left[\prod_{i=1}^n f_i(X_i)\right] = \prod_{i=1}^n \EE[f_i(X_i)], \]
and the expectation on the left-hand side is well defined. The same statement holds for complex-valued functions $f_i$ under the same integrability assumption.

\item As a special case of (1), if $X_1,\dots,X_n$ are independent random variables in $\L^1$, then $X_1\cdots X_n\in \L^1$ and
\[ \EE[X_1\cdots X_n] = \prod_{i=1}^n \EE[X_i]. \]
In general, the product of random variables in $\L^1$ need not belong to $\L^1$ without independence.

\item If $X_1$ and $X_2$ are independent real-valued random variables in $L^2$, then
\[ \operatorname{cov}(X_1,X_2)=0. \]

\item Suppose that, for every $i\in\{1,\dots,n\}$, the random variable $X_i$ admits a density $p_i$ with respect to Lebesgue measure and that $X_1,\dots,X_n$ are independent. Then the random vector $(X_1,\dots,X_n)$ admits a density $p$ given by
\[ p(x_1,\dots,x_n)=\prod_{i=1}^n p_i(x_i). \]

\item Conversely, suppose that $(X_1,\dots,X_n)$ admits a density $p$ of the form
\[ p(x_1,\dots,x_n)=\prod_{i=1}^n q_i(x_i), \]
where $q_1,\dots,q_n$ are nonnegative measurable functions on $\RR$. Then $X_1,\dots,X_n$ are independent, and for every $i\in\{1,\dots,n\}$, $X_i$ admits a density $p_i$ of the form $p_i=C_i q_i$, where $C_i>0$ is a normalization constant.
\end{enumerate}
}{
\begin{enumerate}
    \item  Using the Fubini-Tonelli theorem \ref{thm:Fubini_Tonelli} for the case of two independent variables : 
    \begin{align*}
    \EE\bigl(f_1(X_1)f_2(X_2)\bigr)
            &= \int_{\RR^2} f_1(x_1)f_2(x_2), d\PP_{(X_1,X_2)}(x_1,x_2) \\
            &= \int_{\RR^2} f_1(x_1)f_2(x_2), d(\PP_{X_1}\otimes \PP_{X_2})(x_1,x_2) \\
            &= \int_{\RR} \int_{\RR} f_1(x_1)f_2(x_2), d\PP_{X_2}(x_2), d\PP_{X_1}(x_1) \\
            &= \int_{\RR} f_1(x_1), d\PP_{X_1}(x_1)
            \int_{\RR} f_2(x_2), d\PP_{X_2}(x_2) \\
            &= \EE\bigl(f_1(X_1)\bigr)\times \EE\bigl(f_2(X_2)\bigr).
    \end{align*}
    The general case can be achived by a simple reccurence. 
    \item Similar to (1), but using the Fubini-Lebesgue theorem \ref{thm:Fubini_Lebesgue} and ensuring the $f_i(X_i)\in \L^1$ which is the condition :
    \[ \EE[|f_i(X_i)|]<\infty \quad\text{for every } i\in\{1,\dots,n\}. \]
    \item Clear conclusion from (1). 
    \item This follows from the preceding remarks since :
    \[ \operatorname{Cov}(X_1,X_2) = \EE(X_1X_2) - \EE(X_1)\EE(X_2) = 0 \] 
    \item This is an immediate consequence of the caracterisation of independece $\PP_{(X_1,\ldots,X_n)} = \PP_{X_1} \otimes \cdots \otimes \PP_{X_n}$. Indeed, by the Fubini theorem, under the assumption
    \[ \PP_{X_i}(dx_i) = p_i(x_i)\,dx_i \quad \text{for } i \in \{1,\ldots,n\}, \]
    we have
    \[ \PP_{X_1} \otimes \cdots \otimes \PP_{X_n}(dx_1\cdots dx_n) = \left( \prod_{i=1}^n p_i(x_i) \right) dx_1\cdots dx_n. \]

    \item We first observe, again thanks to the Fubini theorem, that
    \[ \prod_{i=1}^n \left( \int q_i(x)\,dx \right) = \int_{\RR^n} p(x_1,\ldots,x_n)\,dx_1\cdots dx_n = 1. \]
    In particular, the quantity
    \[K_i = \int q_i(x)\,dx\]
    is positive and finite for every \( i \in \{1,\ldots,n\} \).
    Then, by Proposition~8.6, the density of \( X_i \) is
    \begin{align*}
    p_i(x_i)
        &= \int_{\RR^{n-1}} p(x_1,\ldots,x_n)\, dx_1\cdots dx_{i-1}\,dx_{i+1}\cdots dx_n \\
        &= \left( \prod_{j\neq i} K_j \right) q_i(x_i) = \frac{1}{K_i}\, q_i(x_i).
    \end{align*}
    Thus \( p_i = C_i q_i \) with \( C_i = 1/K_i \).
    This also allows us to rewrite the density of \( (X_1,\ldots,X_n) \) as
    \[ p(x_1,\ldots,x_n) = \prod_{i=1}^n q_i(x_i) = \prod_{i=1}^n p_i(x_i), \]
    and therefore
    \[ \PP_{(X_1,\ldots,X_n)} = \PP_{X_1} \otimes \cdots \otimes \PP_{X_n}, \]
    which shows that \( X_1,\ldots,X_n \) are independent.
\end{enumerate}
}


\defnr{Independence of [infinite] \sfields}{infinite_sigmafield_independence}{
Let $(\G_i)_{i\in I}$ be an arbitrary collection of sub-\sfields of
$\A$. We say that the \sfields $\G_i$, $i\in I$, are
\emph{independent} if, for every finite subset
$\{i_1,\dots,i_p\}\subset I$, the \sfields
$\G_{i_1},\dots,\G_{i_p}$ are independent.

If $(X_i)_{i\in I}$ is an arbitrary collection of random variables, we say that
the random variables $X_i$, $i\in I$, are \emph{independent} if the
\sfields $\sigma(X_i)$, $i\in I$, are independent.

Similarly, for an arbitrary collection $(A_i)_{i\in I}$ of events, we say that
the events $A_i$, $i\in I$, are independent if the \sfields
$\sigma(A_i)$, $i\in I$, are independent.
}


\proppr{sigmafield_independence_separation}{
Let $(X_n)_{n\in\NN}$ be a sequence of independent random variables. Then,
for every $p\in\NN$, the \sfields
\[ 
\G_1=\sigma(X_1,\dots,X_p),
\qquad
\G_2=\sigma(X_{p+1},X_{p+2},\dots)
\]
are independent.
}{
We apply the suffisant condition of \sfields independence given in \ref{thm:suffisant_condition_sigmafield_independence}.
Set
\[ \C_1=\sigma(X_1,\dots,X_p)=\G_1 \]
and
\[ \C_2=\bigcup_{k=p+1}^{\infty} \sigma(X_{p+1},\dots,X_k) \subset \G_2. \]
The class $\C_2$ is closed under finite intersections, contains $\Omega$, and satisfies $\sigma(\C_2)=\G_2$.

By independence of the random variables $(X_n)_{n\in\NN}$, and by the grouping-by-blocks principle, we have for every $C_1\in\C_1$ and $C_2\in\C_2$,
\[ \PP(C_1\cap C_2)=\PP(C_1)\PP(C_2). \]
Therefore, the assumptions of theorem \ref{thm:suffisant_condition_sigmafield_independence} are satisfied, and it follows that the \sfields $\G_1$ and $\G_2$ are independent.
}

\subsection{Sums of Independent Random Variables}
Sums of independent random variables play a central role in probability theory. Their distributional and analytic properties are conveniently described using the notion of convolution of probability measures.
Let $\PP$ and $\QQ$ be two probability measures on $\RR^d$. The \emph{convolution} $\PP*\QQ$ is defined as the pushforward of the product measure $\PP\otimes\QQ$ under the mapping $(x,y)\mapsto x+y$. Equivalently, for every nonnegative measurable function $\varphi$ on $\RR^d$,
\[ \int_{\RR^d} \varphi(z)\,(\PP*\QQ)(dz) = \int_{\RR^d}\!\!\int_{\RR^d} \varphi(x+y)\,d\PP(x)\,d\QQ(y). \]

If $\PP$ and $\QQ$ admit densities $f$ and $g$ with respect to Lebesgue measure, then $\PP*\QQ$ admits the density $f*g$ given by
\[ (f*g)(z)=\int_{\RR^d} f(x)\,g(z-x)\,dx. \]

\proppr{convolution_independence}{
Let $X$ and $Y$ be two independent random variables with values in $\RR^d$.
\begin{enumerate}
\item The law of $X+Y$ is the convolution of the laws of $X$ and $Y$:
\[ \PP_{X+Y}=\PP_X*\PP_Y. \]
In particular, if $X$ and $Y$ admit densities $p_X$ and $p_Y$, then $X+Y$ admits the density $p_X*p_Y$.

\item The characteristic function of $X+Y$ satisfies
\[ \Phi_{X+Y}(\xi)=\Phi_X(\xi)\,\Phi_Y(\xi), \qquad \xi\in\RR^d. \]

\item If $\EE[\|X\|^2]<\infty$ and $\EE[\|Y\|^2]<\infty$, then the covariance matrices satisfy
\[ K_{X+Y}=K_X+K_Y. \]
In particular, if $d=1$ and $X,Y\in L^2$, then
\[ \operatorname{var}(X+Y)=\operatorname{var}(X)+\operatorname{var}(Y). \]
\end{enumerate}
}{
\begin{enumerate}
\item We know that $\PP_{(X,Y)} = \PP_X \otimes \PP_Y$, and thus, for every nonnegative measurable function $\varphi$ on $\RR^d$,
\[ \EE[\varphi(X+Y)] = \int \! \int \varphi(x+y)\,\PP_X(dx)\,\PP_Y(dy) = \int \varphi(z)\,(\PP_X * \PP_Y)(dz), \]
by the definition of $\PP_X * \PP_Y$. The case with density follows from the remarks before the proposition.

\item This follows from the equality
\[\PP_{X+Y} = \PP_X * \PP_Y\]
and the last observation before the proposition.

\item If $X = (X_1,\ldots,X_d)$ and $Y = (Y_1,\ldots,Y_d)$, proposition \ref{prop:Random_variables_Independence_proprieties} implies that
\[ \operatorname{cov}(X_i,Y_j) = 0 \quad \text{for every } i,j \in \{1,\ldots,d\}. \]
Consequently, using bilinearity,
\[ \operatorname{cov}(X_i+Y_i,\,X_j+Y_j) = \operatorname{cov}(X_i,X_j) + \operatorname{cov}(Y_i,Y_j), \]
which gives
\[ K_{X+Y} = K_X + K_Y. \]
\end{enumerate}

}

\subsection{The Borel--Cantelli Lemma}
Let $(A_n)_{n\in\NN}$ be a sequence of events. Recall that
\[ \limsup_{\goto{n}} A_n = \bigcap_{n=1}^{\infty}\bigcup_{k=n}^{\infty} A_k. \]
The event $\limsup_{\goto{n}} A_n$ consists of those $\omega\in\Omega$ that belong to infinitely many of the events $A_n$. Note that Lemma~1.7 implies
\[ \PP\bigl(\limsup_{\goto{n}} A_n\bigr) \ge \limsup_{\goto{n}} \PP(A_n). \]

\thmpr{Borel--Cantelli Lemma}{Borel_Cantelli}{
Let $(A_n)_{n\in\NN}$ be a sequence of events.
\begin{enumerate}
\item If $\sum_{n\in\NN} \PP(A_n)<\infty,$ then
\[ \PP\bigl(\limsup_{\goto{n}} A_n\bigr)=0. \]
Equivalently,
\[ \bigl\{n\in\NN:\ \omega\in A_n\bigr\} \quad\text{is finite almost surely.} \]

\item
If $\sum_{n\in\NN} \PP(A_n)=\infty$ and the events $(A_n)_{n\in\NN}$ are independent, then
\[ \PP\bigl(\limsup_{\goto{n}} A_n\bigr)=1. \]
Equivalently,
\[ \bigl\{n\in\NN:\ \omega\in A_n\bigr\} \quad\text{is infinite almost surely.} \]
\end{enumerate}
}{
\begin{enumerate}
\item If $\sum_{n\in\NN} \PP(A_n)<\infty$, then
\[ \EE\!\left[\sum_{n\in\NN} \1_{A_n}\right] = \sum_{n\in\NN} \PP(A_n) < \infty. \]
It follows that
\[ \sum_{n\in\NN} \1_{A_n}<\infty \quad\text{almost surely}, \]
which means that $\omega$ belongs to only finitely many of the events $A_n$. Hence,
\[ \PP\bigl(\limsup_{\goto{n}} A_n\bigr)=0. \]

\item
Assume now that $\sum_{n\in\NN} \PP(A_n)=\infty$ and that the events $A_n$ are independent. Fix $n_0\in\NN$. For $n\ge n_0$, we have
\[ \PP\!\left(\bigcap_{k=n_0}^{n} A_k^c\right) = \prod_{k=n_0}^{n} \PP(A_k^c) = \prod_{k=n_0}^{n} \bigl(1-\PP(A_k)\bigr). \]
Since the series $\sum_{k\in\NN} \PP(A_k)$ diverges, the right-hand side converges to $0$ as $\goto{n}$. Consequently,
\[ \PP\!\left(\bigcap_{k=n_0}^{\infty} A_k^c\right)=0. \]
As this holds for every $n_0\in\NN$, we obtain
\[ \PP\!\left(\bigcup_{n_0=1}^{\infty} \bigcap_{k=n_0}^{\infty} A_k^c\right)=0. \]
Taking complements yields
\[ \PP\!\left(\bigcap_{n_0=1}^{\infty} \bigcup_{k=n_0}^{\infty} A_k\right)=1, \]
that is,
\[ \PP\bigl(\limsup_{\goto{n}} A_n\bigr)=1. \]
\end{enumerate}
}

\section{Conditioning}

\subsection{Definition of conditional expectaiton} 
\subsubsection{Discret case}
As in the previous chapters, we consider a probability space $(\probspace)$. If $B \in \A$ is an event of positive probability,
we can define a new probability measure on $(\measspace)$ by setting, for every $A \in \A$, \[ \PP(A \mid B) := \frac{\PP(A \cap B)}{\PP(B)}. \]

This probability measure $A \mapsto \PP(A \mid B)$ is called \textbf{the conditional probability given $B$}. Similarly, for every nonnegative random variable $X$, or for $X \in L^{1}(\probspace)$,
the \textbf{conditional expectation of $X$ given $B$} is defined by
\[ \EE[X \mid B] := \frac{\EE[X \1_B]}{\PP(B)}. \]

One verifies that $\EE[X \mid B]$ is also the expected value of $X$ under $\PP(\,\cdot \mid B)$.
We can interpret $\EE[X \mid B]$ as the mean value of $X$ knowing that the event $B$ occurs.

Let us now define the conditional expectation knowing a discrete random variable.
We consider a random variable $Y$ taking values in a countable space $E$ equipped as usual
with the $\sigma$-field of all its subsets.
Let
\[ E' = \{ y \in E : \PP(Y = y) > 0 \}. \]
If $X \in L^{1}(\Omega, \A, \PP)$, we may consider, for every $y \in E'$,
\[ \EE[X \mid Y = y] = \frac{\EE[X \1_{\{Y=y\}}]}{\PP(Y = y)}, \]
as a special case of the conditional expectation $\EE[X \mid B]$ when $\PP(B) > 0$.

The conditional expectation of $X$ knowing $Y$ is the real random variable
\[\EE[X \mid Y] = \varphi(Y),\]
where the function $\varphi : E \to \RR$ is given by
\[\varphi(y) =
\begin{cases}
\EE[X \mid Y = y], & \text{if } y \in E',\\
0, & \text{if } y \in E \setminus E'.
\end{cases} \]

\subsubsection{General case}
Now we turn to the definition of conditional expectation with respect to a sub-\sfield. For that purpose we need the next theorem that provides the existance and uniqueness of conditional expectation.
\thmpr{Existance of Conditional Expectation}{Conditional_expectation_existence}{
Let $\B$ be a sub-$\sigma$-field of $\A$, and let $X \in L^1(\Omega,\A,\PP)$. There exists a unique element of $L^1(\Omega,\B,\PP)$, which is denoted by $\EE[X \mid \B]$, such that
\begin{equation} \label{eq:conditional_expectation_characteristic_property1}\forall B \in \B, \quad \EE[X \1_B] = \EE[\EE[X \mid \B] \1_B]. \tag{i} 
\end{equation}

We have more generally, for every bounded $\B$-measurable real random variable $Z$,
\begin{equation} \label{eq:conditional_expectation_characteristic_property2}
\EE[XZ] = \EE[\EE[X \mid \B] Z]. \tag{ii}
\end{equation}
If $X \ge 0$, we have $\EE[X \mid \B] \ge 0$ a.s.
}{
Let us start by proving uniqueness. Let $X'$ and $X''$ be two random variables in $L^1(\Omega,\B,\PP)$ such that
\[ \forall B \in \B, \quad \EE[X' \1_B] = \EE[X \1_B] = \EE[X'' \1_B]. \]
Taking $B = \{X' > X''\}$ (which is in $\B$ since both $X'$ and $X''$ are $\B$-measurable), we get
\[ \EE[(X' - X'') \1_{\{X' > X''\}}] = 0, \]
which implies that $X' \le X''$ a.s., and we have similarly $X' \ge X''$ a.s. Thus $X' = X''$ a.s., which means that $X'$ and $X''$ are equal as elements of $L^1(\Omega,\B,\PP)$.

Let us now turn to existence. We first assume that $X \ge 0$ and $\EE(X)\neq 0$, then we let $\QQ$ be the finite measure on $(\Omega,\B)$ defined by
\[ \forall B \in \B, \quad \QQ(B) := \frac{1}{\EE(X)}\EE[X \1_B]. \]

Let us emphasize that we define $\QQ(B)$ only for $B \in \B$. We may also view $\PP$ as a probability measure on $(\Omega,\B)$, by restricting the mapping $B \mapsto \PP(B)$ to $B \in \B$, and it is immediate that $\QQ \ll \PP$. The Radon--Nikodym theorem \ref{thm:RadonNikodym} applied to the probability measures $\PP$ and $\QQ$ on the measurable space $(\Omega,\B)$ yields the existence of a nonnegative $\B$-measurable random variable $\tilde X$ such that
\[ \forall B \in \B, \quad \EE[X \1_B] = \QQ(B) = \EE[\tilde X \1_B]. \]
Taking $B = \Omega$, we have $\EE[\tilde X] = \EE[X] < \infty$, and thus $\tilde X \in L^1(\Omega,\B,\PP)$. The random variable $\EE[X \mid \B] = \tilde X$ satisfies (1). (for the case of $\EE(X)=0$ we have $X=0$ a.s. and we can take $\tilde{X}=0$)

When $X$ is of arbitrary sign, we just have to take
\[ \EE[X \mid \B] = \EE[X^+ \mid \B] - \EE[X^- \mid \B], \]
and it is clear that (1) also holds in that case.

Finally, to see that (1) implies (2), we rely on the usual measure-theoretic arguments. (2) follows from (1) when $Z$ is a simple random variable (taking only finitely many values), and in the general case lemma \ref{lem:approximation_simple_random_variables} allows us to write $Z$ as the pointwise limit of a sequence $(Z_n)_{n \in \NN}$ of simple $\B$-measurable random variables that are uniformly bounded by the same constant $K$ (such that $|Z| \le K$) and the dominated convergence theorem yields the desired result.
}

The crucial point is the fact that $\EE[X \mid \B]$ is $\B$-measurable. Either of properties (1) and (2) characterizes the conditional expectation $\EE[X \mid \B]$ among random variables of $L^1(\Omega,\B,\PP)$. In what follows, we will refer to (1) or (2) as the \emph{characteristic property} of $\EE[X \mid \B]$.

In the special case where $\B = \sigma(Y)$ is generated by a random variable $Y$, we will write indifferently
\[ \EE[X \mid \B] = \EE[X \mid \sigma(Y)] = \EE[X \mid Y]. \]


We now turn to the definition of $\EE[X \mid \B]$ for a nonnegative random variable $X$ which allows $X$ to take the value $+\infty$.

\thmpr{Existence of Conditional Expectation for nonnegative random variables}{Conditional_expectation_existence_nonnegative}{
Let $X$ be a random variable with values in $[0,\infty]$. There exists a $\B$-measurable random variable with values in $[0,\infty]$, which is denoted by $\EE[X \mid \B]$ and is such that, for every nonnegative $\B$-measurable random variable $Z$,
\[ \EE[XZ] = \EE[\EE[X \mid \B] Z]. \tag{iii} \]
Furthermore $\EE[X \mid \B]$ is unique up to a $\B$-measurable set of probability zero.
}{
We define $\EE[X \mid \B]$ by setting
\[ \EE[X \mid \B] := \lim_{n \uparrow \infty} \EE[X \wedge n \mid \B] \quad \text{a.s.} \]
This definition makes sense because, for every $n \in \NN$, $X \wedge n$ is bounded hence integrable and thus $\EE[X \wedge n \mid \B]$ is well defined by the previous section. Furthermore, the fact that the sequence $(\EE[X \wedge n \mid \B])_{n \in \NN}$ is (a.s.) increasing follows from the fact that the conditional expectation of a.s. non-negative random variable is also an a.s. non-negative random variable. Then, if $Z$ is nonnegative and $\B$-measurable, the monotone convergence theorem implies that
\[ \EE[\EE[X \mid \B] Z] = \lim_{n \to \infty} \EE[\EE[X \wedge n \mid \B] Z \1_{\{Z < \infty\}}] = \lim_{n \to \infty} \EE[(X \wedge n) Z \1_{\{Z < \infty\}}] = \EE[XZ]. \]

It remains to establish uniqueness. Let $X'$ and $X''$ be two $\B$-measurable random variables with values in $[0,\infty]$, such that
\[ \EE[X'Z] = \EE[X''Z] \]
for every nonnegative $\B$-measurable random variable $Z$. Let us fix two nonnegative rationals $a < b$, and take
\[ Z = \1_{\{X' \le a < b \le X''\}}. \]

It follows that
\[ a \, \PP(X' \le a < b \le X'') \le \EE[\1_{\{X' \le a < b \le X''\}} X'] = \EE[\1_{\{X' \le a < b \le X''\}} X''] \ge b \, \PP(X' \le a < b \le X''), \]
which is only possible if $\PP(X' \le a < b \le X'') = 0$. Hence,
\[ \PP\!\left( \bigcup_{\substack{a,b \in \QQ_+ \\ a < b}} \{ X' \le a < b \le X'' \} \right) = 0, \]
which implies $X' \ge X''$ a.s., and interchanging the roles of $X'$ and $X''$ also gives $X'' \ge X'$ a.s.
}

\textbf{Remark:}
The uniqueness part of the theorem means that we should view $\EE[X \mid \B]$ as an equivalence class of $\B$-measurable random variables that are equal a.s. But of course it will be more convenient to speak about "the random variable" $\EE[X \mid \B]$.

In the case where $X$ is also integrable, the definition of the preceding theorem is consistent with that of Theorem \ref{thm:Conditional_expectation_existence}. Indeed (3) with $Z = 1$ shows that $\EE[\EE[X \mid \B]] = \EE[X]$, and in particular $\EE[X \mid \B] < \infty$ a.s.\ so that we can take a representative with values in $[0,\infty)$ and $\EE[X \mid \B] \in L^1$. We then just have to note that (1) (in the theorem we are citing) is the special case of (3) where $Z = \1_B$, $B \in \B$.

Similarly as in the case of integrable variables, we will refer to (3) (or to its variant when $Z = \1_B$, $B \in \B$) as the characteristic property of $\EE[X \mid \B]$.

\subsection{Properties of conditional Expectations}
We now list the properties of conditional expectaitons : 
\proppr{conditional_expectation_proprieties}{
\begin{enumerate}
\item  For a $X,Y \in L^1(\Omega,\A,\PP)$)  and $\alpha,\beta\in\RR$ we have : 
\[ \EE(\alpha X + \beta Y \mid \B ) = \alpha \EE(X\mid \B) + \beta\EE(Y\mid \B) \] 
\item  For a $X,Y$ non negative random variables and $\alpha,\beta\in\RR^+$ we have : 
\[ \EE(\alpha X + \beta Y \mid \B ) = \alpha \EE(X\mid \B) + \beta\EE(Y\mid \B) \] 
\item If $X \in L^1(\Omega,\A,\PP)$) (or if $X$ is just non negative) and $X$ is $\B$-measurable, then :
\[\EE[X \mid \B] = X\]
\item If $X \in L^1(\Omega,\A,\PP)$) (or if $X$ is just non negative), then :
\[\EE[\EE[X \mid \B]] = \EE[X]\]
\item If $X, X' \in L^1(\Omega,\A,\PP)$) and $X \ge X'$, then : 
\[ \EE[X \mid \B] \ge \EE[X' \mid \B] \quad \text{a.s} \]
\item If $X \in L^1(\Omega,\A,\PP)$), then :
\[ |\EE[X \mid \B]| \le \EE[|X| \mid \B] \quad \text{a.s.}\] 
and, consequently,
\[ \EE[|\EE[X \mid \B]|] \le \EE[|X|]. \]
Therefore the mapping $X \mapsto \EE[X \mid \B]$ is a contraction of $L^1(\Omega,\A,\PP)$).
\end{enumerate}
}{
    The proofs for these properties are left exercise for the reader.
}
Now the properties allowing intereaction between limits and conditional expectations : 
\proppr{conditional_expectation_limits_properties}{
\begin{enumerate}
\item If $(X_n)_{n \in \NN}$ is an increasing sequence of nonnegative random variables, and $X = \lim \uparrow X_n$, then
\[ \EE[X \mid \B] = \lim_{n \to \infty} \uparrow \, \EE[X_n \mid \B], \quad \text{a.s.} \]

As consequence : if $(Y_n)_{n \in \NN}$ is a sequence of nonnegative random variables, we have
\[ \EE\!\left[ \sum_{n \in \NN} Y_n \,\middle|\, \B \right] = \sum_{n \in \NN} \EE[Y_n \mid \B]. \]

\item If $(X_n)_{n \in \NN}$ is any sequence of nonnegative random variables, then
\[ \EE[\liminf X_n \mid \B] \le \liminf \EE[X_n \mid \B], \quad \text{a.s.} \]

\item Let $(X_n)_{n \in \NN}$ be a sequence of integrable random variables that converges a.s.\ to $X$. Assume that there exists a nonnegative random variable $Z$ such that $|X_n| \le Z$ a.s.\ for every $n \in \NN$, and $\EE[Z] < \infty$. Then,
\[ \EE[X \mid \B] = \lim_{n \to \infty} \EE[X_n \mid \B], \quad \text{a.s.}  \]
The convergence is also in $\L^1$ that is : 
\[ \EE\left[ \big| \EE[X \mid \B] - \EE[X_n \mid \B] \big| \right] \longrightarrow 0 \quad \textit{as } n\rightarrow +\infty  \]
\end{enumerate}
}{
\begin{enumerate}
\item It follows from last property that we have $\EE[X \mid \B] \ge \EE[Y \mid \B]$ if $X \ge Y \ge 0$. Under the assumptions of (1), we can therefore set
\[ X' = \lim \uparrow \EE[X_n \mid \B], \]
which is a $\B$-measurable random variable with values in $[0,\infty)$. Then, for every nonnegative $\B$-measurable random variable $Z$, the monotone convergence theorem gives
\[ \EE[ZX'] = \lim \uparrow \EE[Z \EE[X_n \mid \B]] = \lim \uparrow \EE[Z X_n] = \EE[ZX], \]
which implies $X' = \EE[X \mid \B]$ thanks to the characteristic property (3) from theorem \ref{thm:Conditional_expectation_existence}. The second assertion in (1) follows by applying the first one to $X_n = Y_1 + \cdots + Y_n$.

\item Using (1), we have
\[ \EE[\liminf X_n \mid \B] = \EE\!\left[ \lim_{k \uparrow \infty} \left( \inf_{n \ge k} X_n \right) \,\middle|\, \B \right] = \lim_{k \uparrow \infty} \EE\!\left[ \inf_{n \ge k} X_n \,\middle|\, \B \right]. \]

\item It suffices to apply (2) twice:
\[ \EE[Z - X \mid \B] = \EE[\liminf (Z - X_n) \mid \B] \le \EE[Z \mid \B] - \limsup \EE[X_n \mid \B], \]
\[ \EE[Z + X \mid \B] = \EE[\liminf (Z + X_n) \mid \B] \le \EE[Z \mid \B] + \liminf \EE[X_n \mid \B], \]
which leads to
\[ \EE[X \mid \B] \le \liminf \EE[X_n \mid \B] \le \limsup \EE[X_n \mid \B] \le \EE[X \mid \B], \]
giving the desired almost sure convergence. The convergence in $L^1$ is now a consequence of the dominated convergence theorem, since we have
\[ |\EE[X_n \mid \B]| \le \EE[|X_n| \mid \B] \le \EE[Z \mid \B] \]
and $\EE[\EE[Z \mid \B]] = \EE[Z] < \infty$.
\end{enumerate}
}

\thmpr{Jensen Inequality conditional Expectation}{conditional_Jensen_inequality}{
If $f : \RR \to \RR_+$ is convex, and if $X \in L^1$, then
\[ \EE[f(X) \mid \B] \ge f(\EE[X \mid \B]). \]
}{
Set
\[ E_f = \{ (a,b) \in \RR^2 : \forall x \in \RR, \ f(x) \ge ax + b \}. \]
Then,
\[ \forall x \in \RR, \quad f(x) = \sup_{(a,b) \in E_f} (ax + b) = \sup_{(a,b) \in E_f \cap \QQ^2} (ax + b). \]

We can take advantage of the fact that $\QQ^2$ is countable to discard a countable collection of sets of probability zero and to get that, a.s.,
\[ \EE[f(X) \mid \B] = \EE\!\left[ \sup_{(a,b) \in E_f \cap \QQ^2} (aX + b) \,\middle|\, \B \right] \ge \sup_{(a,b) \in E_f \cap \QQ^2} \EE[aX + b \mid \B] = f(\EE[X \mid \B]). \quad  \]
}
As a consequence of Jensen's inequality, we obtain that, for every $p \ge 1$, the mapping
$X \mapsto \EE[X \mid \B]$ maps $L^p(\Omega,\A,\PP)$ into itself, and is a contraction of $L^p(\Omega,\A,\PP)$. Indeed, for every $X \in L^p(\Omega,\A,\PP)$, we have using property (d) in Section~11.2.1,
\[ \EE[|\EE[X \mid \B]|^p] \le \EE[\EE[|X| \mid \B]^p] \le \EE[\EE[|X|^p \mid \B]] = \EE[|X|^p]. \]

This contraction property also holds for $p = \infty$.

\thmp{The orthogonal projection property}{
If $X \in L^2(\Omega,\A,\PP)$, then $\EE[X \mid \B]$ is the orthogonal projection of $X$ on $L^2(\Omega,\B,\PP)$.
}{
Jensen's inequality shows that $\EE[X \mid \B]^2 \le \EE[X^2 \mid \B]$, a.s. This implies that
\[
\EE[\EE[X \mid \B]^2] \le \EE[\EE[X^2 \mid \B]] = \EE[X^2] < \infty,
\]
and thus the random variable $\EE[X \mid \B]$ belongs to $L^2(\Omega,\B,\PP)$.

On the other hand, for every bounded $\B$-measurable random variable $Z$,
\[
\EE[Z(X - \EE[X \mid \B])] = \EE[ZX] - \EE[Z\EE[X \mid \B]] = 0,
\]
by the characteristic property (11.2). Hence $X - \EE[X \mid \B]$ is orthogonal to the space of bounded $\B$-measurable random variables, and the latter space is dense in $L^2(\Omega,\B,\PP)$ (for instance by Theorem~4.8). It follows that $X - \EE[X \mid \B]$ is orthogonal to $L^2(\Omega,\B,\PP)$, which gives the desired result.
}

\thmpr{Pull-out property of conditional expectation}{Pull_out_property}{
Let $X$ and $Y$ be real random variables and assume that $Y$ is $\B$-measurable. Then
\[ \EE[YX \mid \B] = Y\,\EE[X \mid \B], \]
provided $X$ and $Y$ are both nonnegative, or $X$ and $YX$ are both integrable.
}{
Suppose that $X \ge 0$ and $Y \ge 0$. Then, for every nonnegative $\B$-measurable random variable $Z$, we have
\[ \EE[Z(Y\EE[X \mid \B])] = \EE[(ZY)\EE[X \mid \B]] = \EE[ZYX], \]
using \textit{characteristic property} (3) from theorem \ref{thm:Conditional_expectation_existence} and the fact that $ZY$ is $\B$-measurable. Since $Y\EE[X \mid \B]$ is a nonnegative $\B$-measurable random variable, the characteristic property from theorem \ref{thm:Conditional_expectation_existence} implies that
\[ Y\EE[X \mid \B] = \EE[YX \mid \B]. \]

In the case when $X$ and $YX$ are integrable, we get the desired result by decomposing
\[ X = X^+ - X^-, \qquad Y = Y^+ - Y^- . \]
}

\thmpr{Nested $\sigma$-Fields}{conditional_nested_sigma_fields}{
Let $\B_1$ and $\B_2$ be two sub-$\sigma$-fields of $\A$ such that $\B_1 \subset \B_2$. Then, for every nonnegative (or integrable) random variable $X$, we have
\[ \EE[\EE[X \mid \B_2] \mid \B_1] = \EE[X \mid \B_1]. \]
}{
Consider the case where $X \ge 0$. Let $Z$ be a nonnegative $\B_1$-measurable random variable. Then, since $Z$ is also $\B_2$-measurable,
\[ \EE[Z\,\EE[\EE[X \mid \B_2] \mid \B_1]] = \EE[Z\,\EE[X \mid \B_2]] = \EE[ZX]. \]
Hence $\EE[\EE[X \mid \B_2] \mid \B_1]$ satisfies the characteristic property of $\EE[X \mid \B_1]$, and the desired result follows.
}
\textbf{Remark:} We have also $\EE[\EE[X \mid \B_1] \mid \B_2] = \EE[X \mid \B_1]$ under the same assumptions, but this is trivial since $\EE[X \mid \B_1]$ is $\B_2$-measurable.

\thmp{Independence and conditional expectaiton}{
Two sub-$\sigma$-fields $\B_1$ and $\B_2$ of $\A$ are independent if and only if, for every $B \in \B_2$, we have
\[ \EE[\1_B \mid \B_1] = \PP(B). \]
Furthermore, if $\B_1$ and $\B_2$ are independent, we have also $\EE[X \mid \B_1] = \EE[X]$ for every nonnegative $\B_2$-measurable random variable $X$ and for every $X \in L^1(\Omega,\B_2,\PP)$.
}{
Left exercise for the reader.
}

\thmpr{Factorization under independence}{factorization_under_independence_property}{
Let $(E,\E)$ and $(F,\F)$ be two measurable spaces, and let $X$ and $Y$ be random variables taking values in $E$ and $F$ respectively. Assume that $X$ is independent of $\B$ and that $Y$ is $\B$-measurable. Then, for every $\E \otimes \F$-measurable function $g : E \times F \to \RR_+$,
\[ \EE[g(X,Y) \mid \B] = \int g(x,Y)\,\PP_X(dx), \]
where we recall that $\PP_X$ denotes the law of $X$. The right-hand side is the composition of the random variable $Y$ with the mapping $\Psi : F \to \RR_+$ defined by
\[ \Psi(y) = \int g(x,y)\,\PP_X(dx). \]
}{
We have to prove that, for any nonnegative $\B$-measurable random variable $Z$, we have
\[ \EE[g(X,Y)Z] = \EE[\Psi(Y)Z]. \]

Write $\PP_{X,Y,Z}$ for the law of the triple $(X,Y,Z)$, which is a probability measure on $E \times F \times \RR_+$. Since $X$ is independent of $\B$, $X$ is independent of the pair $(Y,Z)$, and therefore
\[ \PP_{X,Y,Z} = \PP_X \otimes \PP_{Y,Z}. \]
Then, using the Fubini theorem, we have
\begin{align*}
\EE[g(X,Y)Z]
&= \int g(x,y)z \, \PP_{X,Y,Z}(dx\,dy\,dz) \\
&= \int g(x,y)z \, \PP_X(dx)\,\PP_{Y,Z}(dy\,dz) \\
&= \int_{F \times \RR_+} z \left( \int_E g(x,y)\,\PP_X(dx) \right) \PP_{Y,Z}(dy\,dz) \\
&= \int_{F \times \RR_+} z\,\Psi(y)\,\PP_{Y,Z}(dy\,dz) \\
&= \EE[\Psi(Y)Z].
\end{align*}
which was the desired result. 
}
\textbf{Remark:}
The mapping $\Psi$ is measurable thanks to the Fubini theorem. The theorem can be explained informally as follows. If we condition with respect to the $\sigma$-field $\B$, the random variable $Y$, which is $\B$-measurable, behaves like a constant, but on the other hand $\B$ gives no information on the random variable $X$. Thus the best approximation of $g(X,Y)$ knowing $\B$ is obtained by integrating $g(\cdot,Y)$ with respect to the law of $X$.

\thmpr{Invariance under independent information}{invariance_under_independent_information_property}{
Let $Z$ be a random variable in $L^1$, and let $\H_1$ and $\H_2$ be two sub-$\sigma$-fields of $\A$. Assume that $\H_2$ is independent of $\sigma(Z) \vee \H_1$. Then,
\[\EE[Z \mid \H_1 \vee \H_2] = \EE[Z \mid \H_1].\]
}{
It suffices to prove that the equality
\[ \EE[\1_A Z] = \EE[\1_A \EE[Z \mid \H_1]]  \]
holds.
}

\subsection{Transition Probabilities and Conditional Distributions}
\defn{Transition probability}{
Let $(E,\E)$ and $(F,\F)$ be two measurable spaces. A transition probability (or transition kernel) from $E$ into $F$ is a mapping
\[ \nu : E \times \F \longrightarrow [0,1] \]
which satisfies the following two properties:
\begin{enumerate}
\item for every $x \in E$, $A \mapsto \nu(x,A)$ is a probability measure on $(F,\F)$;
\item for every $A \in \F$, the function $x \mapsto \nu(x,A)$ is $\E$-measurable.
\end{enumerate}
}

At an intuitive level, each time one fixes a "starting point" $x \in E$, the probability measure $\nu(x,\cdot)$ gives a way to choose an "arrival point" $y \in F$ in a random manner (but with a law depending on the starting point $x$). This notion plays a fundamental role in the theory of Markov chains.

\propp{
Let $\nu$ be a transition probability from $E$ into $F$.
\begin{enumerate}
\item If $h$ is a nonnegative (resp.\ bounded) measurable function on $(F,\F)$, then
\[ \varphi(x) = \int h(y)\,\nu(x,dy), \qquad x \in E, \]
is a nonnegative (resp.\ bounded) measurable function on $E$.
\item If $\gamma$ is a probability measure on $(E,\E)$, then
\[\mu(A) = \int \gamma(dx)\,\nu(x,A), \qquad A \in \F,\]
is a probability measure on $(F,\F)$.
\end{enumerate}
}{
We omit the easy proof of the proposition. When treating the nonnegative case of (1), we first consider simple functions, and then use an increasing passage to the limit.  }

We now come to the relation between the notion of a transition probability and conditional expectations.

\defn{The conditional distribution}{
Let $X$ and $Y$ be two random variables taking values in $(E,\E)$ and $(F,\F)$ respectively. Any transition probability $\nu$ from $E$ into $F$ such that, for every nonnegative measurable function $h$ on $F$,
\[ \EE[h(Y) \mid X] = \int h(y)\,\nu(X,dy), \qquad \text{a.s.}, \]
is called a conditional distribution of $Y$ knowing $X$.
}

\textbf{Remark}
The random variable $\int h(y)\,\nu(X,dy)$ is the composition of $X$ with the function
\[ x \mapsto \int h(y)\,\nu(x,dy), \]
which is measurable by last Proposition. In particular, this random variable is a function of $X$, as $\EE[h(Y) \mid X]$ should be.

By definition, if $\nu$ is a conditional distribution of $Y$ knowing $X$, we have, for every $A \in \F$,
\[ \PP(Y \in A \mid X) = \nu(X,A), \qquad \text{a.s.} \]

It is tempting to replace this equality of random variables by an equality of real numbers,
\[ \PP(Y \in A \mid X = x) = \nu(x,A), \]
for every $x \in E$. Although it gives the intuition behind the notion of conditional distribution, this last equality has no mathematical meaning in general, because one has typically $\PP(X = x) = 0$ for every $x \in E$, which makes it impossible to define the conditional probability given the event $\{X = x\}$. The only rigorous formulation is thus the first equality $\PP(Y \in A \mid X) = \nu(X,A)$.

Let us briefly discuss the uniqueness of the conditional distribution of $Y$ knowing $X$. If $\nu$ and $\nu'$ are two such conditional distributions, we have, for every $A \in \F$,
\[ \nu(X,A) = \PP(Y \in A \mid X) = \nu'(X,A), \qquad \text{a.s.} \]
which is equivalent to saying that, for every $A \in \F$,
\[\nu(x,A) = \nu'(x,A), \qquad \PP_X(dx)\text{-a.s.}\]

Suppose that the measurable space $(F,\F)$ is such that any probability measure on $(F,\F)$ is characterized by its values on a (given) countable collection of measurable sets (this property holds for $(\RR^d,\B(\RR^d))$, by considering the collection of all rectangles with rational coordinates). Then, we conclude from the preceding display that
\[\nu(x,\cdot) = \nu'(x,\cdot), \qquad \PP_X(dx)\text{-a.s.}\]

So uniqueness holds in this sense (and clearly we cannot expect more). By abuse of language, we will nonetheless speak about \emph{the conditional distribution} of $Y$ knowing $X$.

Let us now consider the problem of the existence of conditional distributions.

\thmp{Existence of conditional distribution}{
Let $X$ and $Y$ be two random variables taking values in $(E,\E)$ and $(F,\F)$, respectively. Suppose that $(F,\F)$ is a complete separable metric space equipped with its Borel $\sigma$-field. Then, there exists a conditional distribution of $Y$ knowing $X$.
}{
We will not prove this theorem, see Theorem 6.3 in \cite{Probability2} for the proof of a more general statement.
}

\subsection{Evaluation of Conditional Expectation}
\subsubsection{Discret Case}
If $X$ is a discrete variable ($E$ is countable), and $Y$ any random variable. then we have the conditional distibution of $Y$ given $X$, denoted by $\nu(x,A)$, defined by :
\[ \nu(x,A) := 
\begin{cases} 
    \frac{1}{\PP(X=x)}\PP(Y \in A, X = x) & \text{if } x \in E' := \{y \in E : \PP(X = y) > 0\}, \\[6pt]
    \delta_{y_0}(A) & \text{if } x \notin E', 
\end{cases} \]
where $y_0$ is an arbitrary fixed point of $F$. Theorefore we can conclude that : 
\[ \EE[h(Y) \mid X] = \varphi(X) \]
where
\[ \varphi(x) = \int h(y)\,\nu(x,dy) = 
\begin{cases} 
    \frac{1}{\PP(X=x)}\EE(h(Y)\1_{\{X=x\}}) & \text{if } x \in E' , \\[6pt]
    y_0 & \text{if } x \notin E', 
\end{cases} \]
Where $y_0$ is an arbitrary fixed point of $F$.

\subsubsection{Continuous Case [Random Variables with a Density]}
Suppose that $X$ and $Y$ take values in $\RR^m$ and in $\RR^n$ respectively, and that the pair $(X,Y)$ has density $p(x,y)$, $(x,y) \in \RR^m \times \RR^n$. The density of $X$ is then
\[q(x) = \int_{\RR^n} p(x,y)\, dy,\]
where we take $q(x) = 0$ if the integral is infinite. We may define the conditional distribution of $Y$ knowing $X$ by :
\[ \nu(x,A) = 
\begin{cases} 
    \frac{1}{q(x)} \int_A p(x,y)\, dy & \text{if } q(x) > 0, \\[6pt]
    \delta_0(A) & \text{if } q(x) = 0. 
\end{cases} 
\]
 Theorefore we can conclude that : 
\[ \EE[h(Y) \mid X] = \varphi(X) \]
where
\[ \varphi(x) = \int h(y)\,\nu(x,dy) = 
\begin{cases}
    \dfrac{1}{q(y)} \displaystyle \int_{\RR^m} h(x)\, p(x,y)\, dx & \text{if } q(y) > 0, \\
    h(0) & \text{if } q(y) = 0.
\end{cases}\]

\section{Gaussian Distribution} \label{sec:gaussian_sec}
\lipsum


\section{Convergence of random variables} 

The study of convergence of random variables is fundamental in probability theory, particularly in the analysis of asymptotic behaviors and limit theorems. Different notions of convergence allow us to formalize how a sequence of random variables approximates another random variable under various probabilistic criteria. These modes of convergence differ in strength and applicability, and understanding their relationships is essential for results such as the law of large numbers and the central limit theorem.

Throughout this section, we consider a probability space $(\probspace)$. Let $(\bX_n)_{n\in\NN}$ be a sequence of random variables that mapsto $\RR^d$ and let $\bX$ be another random variable taking values in $\RR^d$. We introduce below the principal notions of convergence that will be used in the sequel.

\defnr{Almost sure convergence}{as_conv}{
The sequence $(\bX_n)$ is said to converge almost surely to $\bX$, denoted by
\[\bX_n \conv[a.s.]{n} \bX, \]
if
\[\PP\left(\left\{\omega \in \Omega :\lim_{n\to\infty} \bX_n(\omega)=\bX(\omega)\right\}\right)=1.\]}

\proppr{suffisant_condition_for_as_convergence}{
Let $(\bX_n)_{n\geq 1}$ and $\bX$ be integrable random variables. If
\[ \sum_{n=1}^{\infty} \EE{|\bX_n-\bX|}<\infty, \]
then
\[ \bX_n \xrightarrow{\text{a.s.}} \bX. \]
}{
Let $\epsilon>0$ 
By Markov's inequality,
\[ \mathbb P\big(|X_n-X|>\epsilon\big)  \le \frac{\mathbb E|X_n-X|}{\epsilon}. \]
Therefore,
\[
\sum_{n=1}^\infty \mathbb P\big(|X_n-X|>\epsilon\big)  
\le \frac{1}{\varepsilon}\sum_{n=1}^\infty \mathbb E|X_n-X|
<\infty.
\]
By the Borel-Cantelli lemma \ref{thm:Borel_Cantelli} :
\[ \mathbb P\left(\limsup |X_n-X|>\epsilon \right)=0. \]
Now take $\epsilon=1/m$ for $m\in\mathbb N^*$. Then for every $m\ge1$,
\[ \mathbb P\left(\limsup |X_n-X|>\frac{1}{m} \right)=0. \]
Since a countable union of null sets is null,
\[\mathbb P\Big(\bigcup_{m=1}^\infty \{\limsup |X_n-X|>\frac{1}{m} \}\Big)=0.\]
Thus, outside a null set, for every $m\ge1$,
\[ \limsup |X_n-X|\le \frac{1}{m} \quad a.s.\]
Then : 
\[ \limsup |X_n-X| = 0 \quad a.s.\]
This implies that $X_n\to X$ almost surely.}

\defnr{Convergence in probability}{prob_cong}{
We say that the sequence $(\bX_n)_{n\in\NN}$ converges in probability to $\bX$, and write
\[\bX_n \conv[\PP]{n} \bX,\]
if for every $\varepsilon>0$,
\[\lim_{n\to\infty}\PP\big(|\bX_n - \bX|>\varepsilon\big)=0.\]}


\proppr{prob_conv_characterisation}{
We remind that $d$ is the distance on $L^0(\Omega,\A,\PP)$ defined by :
\[ d(\bX,\bY) = \EE[|\bX-\bY| \wedge  1]\] 
We have : 
\[ \bX_n \conv[\PP]{n} \bX \Longleftrightarrow d(\bX_n,\bX) \conv{n} 0\]
}{
Left exercice for the reader. \\  
Hint: use the fact that the space $L^0_{\RR^d}(\Omega,\A,\PP)$ is complete for the distance $d$ by proposition \ref{prop:L0_complete}.
}

\defnr{Convergence in Lp}{Lp_conv} {
Let $p\in[1,\infty)$ and assume that
\[\EE[|\bX|^p]<\infty\quad \text{and} \quad\EE[|\bX_n|^p]<\infty \quad \text{for all } n\in\NN.\]
The sequence $(\bX_n)$ is said to converge to $\bX$ in $L^p$, written
\[\bX_n \conv[\L^p]{n} \bX,\]
if
\[\lim_{n\to\infty}\EE\big[|\bX_n - \bX|^p\big]=0.\]}
This is equivalent to saying that each component sequence of $(\bX_n)_{n \in \NN}$ converges to the corresponding component of $\bX$ in the Banach space $L^p(\Omega,\A,\PP)$. One could also consider the convergence in $L^\infty$, but this convergence is very rarely used in probability theory and we will not discuss it.

\subsection{Relationship between modes of convergences }
\proppr{as_or_Lp_imply_Prob}{
    If : 
    \[ 
        \bX_n \conv[a.s.]{n} \bX, \quad \text{or} \quad \bX_n \conv[\L^p]{n} \bX.
    \]
    Then : 
    \[ 
        \bX_n \conv[\PP]{n} \bX.
    \]
}{
If $\bX_n$ converges almost surely to $\bX$, then by dominated convergence :
\[d(\bX_n,\bX) = \EE\bigl[\,|\bX_n - \bX| \wedge 1\,\bigr] \longrightarrow 0\quad \text{as } n \to \infty,\]

If $\bX_n$ converges to $\bX$ in $L^p$, then
\[d(\bX_n,\bX) \le \EE\bigl[|\bX_n - \bX|\bigr]= \|\bX_n - \bX\|_1\le \|\bX_n - \bX\|_p\longrightarrow 0\quad \text{as } n \to \infty.\]
}

\proppr{prob_imply_as}{
    If : 
    \[ 
        \bX_n \conv[\PP]{n} \bX.
    \]
    Then : 
    \[ 
        \exists, (n_k) \uparrow \infty \quad \text{ such that }\quad  \bX_{n_k} \conv[a.s.]{k} \bX
    \]
}{
The  assertion has been derived in the proof of Proposition \ref{prop:tobeadjusted}.  
}

\proppr{prob_imply_Lp}{
    If : 
    \[ 
    \begin{cases*}
        \bX_n \conv[\PP]{n} \bX. \\
        \exists r \in (1,\infty) \quad \text{ such : }\quad  (\EE[|\bX_n|^r])_{n \in \NN} \quad \text{is bounded}.
    \end{cases*}
    \]
    Then : 
    \[ 
    \begin{cases*}
        \EE[|\bX|^r] < \infty \\ 
        \forall p \in [1,r), \quad \bX_n \conv[\L^p]{n} \bX.
    \end{cases*}
    \]
}{
By assumption, there is a constant $C$ such that $\EE[|\bX_n|^r] \le C$ for every $n$. By the previous proposition we can find a subsequence $(\bX_{n_k})$ that converges almost surely to $\bX$. By Fatou's lemma,

\begin{align*}
\EE[|\bX|^r]
&= \EE\Bigl[\liminf_{k \to \infty} |\bX_{n_k}|^r\Bigr] \\
&\le \liminf_{k \to \infty} \EE[|\bX_{n_k}|^r] \\
&\le C.
\end{align*}

Then, for every $p \in [1,r)$ and every $\varepsilon > 0$,
\begin{align*}
\EE[|\bX_n - \bX|^p] 
    &= \EE\bigl[|\bX_n - \bX|^p \1_{\{|\bX_n - \bX| \le \varepsilon\}}\bigr]  + \EE\bigl[|\bX_n - \bX|^p \1_{\{|\bX_n - \bX| > \varepsilon\}}\bigr] \\
    &\le \varepsilon^p + \EE[|\bX_n - \bX|^r]^{p/r}\PP(|\bX_n - \bX| > \varepsilon)^{1 - p/r} \\
    &\le \varepsilon^p + (2C)^{p/r} \PP(|\bX_n - \bX| > \varepsilon)^{1 - p/r}.
\end{align*}

Since $\bX_n$ converges in probability to $\bX$,
\[\limsup_{n \to \infty} \EE[|\bX_n - \bX|^p] \le \varepsilon^p.\]
Letting $\varepsilon \to 0$ yields $\EE[|\bX_n - \bX|^p] \to 0$.
}

\lempr{Scheffé's lemma}{prob_imply_L1}{
    If : 
    \[ 
    \begin{cases*}
        \bX_n \geq 0 \quad \text{and}\quad  \bX_n \conv[\PP]{n} \bX. \\
        \forall n\in \NN :\quad \bX_n\in L^1 \quad \text{and}\quad \bX \in L^1 \\
        \EE[\bX_n] \conv{n} \EE[\bX]\\
    \end{cases*}
    \]
    Then : 
    \[ 
        \bX_n \conv[\L^1]{n} \bX.
    \]
}{
We write
\[ \EE[|\bX_n - \bX|] = \EE[\bX_n - \bX] + 2\EE[(\bX_n - \bX)^-]. \]
We have $\EE[\bX_n - \bX] \to 0$, and the bound $(\bX_n - \bX)^- \le \bX$ allows us to apply the dominated convergence theorem to obtain $\EE[(\bX_n - \bX)^-] \to 0$ (using that almost sure convergence in the dominated convergence theorem can be replaced by convergence in probability).
}

The next theorem would be seen in the martingale theory (next chapter), but we state it here for the sake of completeness. It gives a equivalent condition for the convergence in $L^1$ under the assumption of convergence in probability. The condition is expressed in terms of uniform integrability, that we state here for information but it's studied deeply in next chapter. \\ 
A collection $(X_i)_{i\in I}$ of random variables in $L^1(\Omega,\A,\PP)$ is said to be \textbf{uniformly integrable} (u.i.) if
\[ \lim_{a\to+\infty} \sup_{i\in I} \EE\big[|X_i|\1_{\{|X_i|>a\}}\big] = 0. \]

\thmpr{Vitali convergence theorem}{prob_imply_L1_vitali}{
If 
\[ \begin{cases*}
    \bX_n \conv[\PP]{n} \bX. \\
    \bX_n \in L^1 \quad \text{for every } n\in\NN.
\end{cases*} \] 
Then the following assertions are equivalent :
\[
\begin{cases}
    1. & \bX_n \conv[L^1]{n} \bX \\ 
    2. & \text{the sequence } (\bX_n)_{n \in \NN} \text{ is u.i.} \\
\end{cases}
\]
}{
    It would be shown in the section of uniform integrability \ref{sec:uniform_integrability} (next chapter).  
}
It's also relevant to see that if there is $Z\in L^1$ such $|\bX_n| \leq Z$ a.s. for every $n\in\NN$, then the sequence $(\bX_n)_{n \in \NN}$ is u.i., which gives a simpler condition to prove $L^1$-convergence. In fact, we have :
\[ \EE[|\bX_n \1_{|\bX_n\geq a}] \leq \EE[Z\1_{Z\geq a}]  \]
Finaly, taking supremum and limits : 
\[ \lim_{\goto{a}} \sup_{n\in\NN} \EE[|\bX_n| \1_{|\bX_n\geq a}] \leq \lim_{\goto{a}} \EE[Z\1_{Z\geq a}] = 0 \quad \text{Because :}\quad Z\1_{Z\geq a} \conv[a.s.]{n} 0\]
 

\begin{figure}[h]
\centering
\[
\begin{tikzcd}[column sep=huge, row sep=large]
|[draw, rounded corners, inner sep=3pt]| {X_n \xrightarrow{L^q} X}
  \arrow[r, "q \ge p \ge 1"', Rightarrow]
&
|[draw, rounded corners, inner sep=3pt]| {X_n \xrightarrow{L^p} X}
  \arrow[r, Rightarrow]
&
|[draw, rounded corners, inner sep=3pt]| {X_n \xrightarrow{L^1} X}
  \arrow[d, Rightarrow]
\\
&&
|[draw, rounded corners, inner sep=3pt]| {X_n \xrightarrow{\PP} X}
  \arrow[r, Rightarrow]
  \arrow[lu, "{L^r-Bounded \quad r>p}", dashed, bend left=35]
  \arrow[u, "{\text{Expectation convergence}}"', dashed, bend right=45]
  \arrow[u, swap, "{\substack{\text{Uniform conv.}\\ \text{or domination}}}"', dashed, bend left=45]
  \arrow[d, "{\text{along a subsequence}}", dashed, bend left=60]
&
|[draw, rounded corners, inner sep=3pt]| {X_n \xrightarrow{\mathcal D} X}
\\
&&
|[draw, rounded corners, inner sep=3pt]| {X_n \xrightarrow{\text{a.s.}} X}
  \arrow[u, Rightarrow]
\end{tikzcd}
\]
\caption{Relations between the different convergences} \label{fig:relationshipsConv}
\end{figure}


Figure \ref{fig:relationshipsConv} illustrates the relations between the different types of convergence that we have introduced. The dashed lines give partial converses (under the stated conditions) that correspond to results stated in Propositions \ref{prop:prob_imply_as}, \ref{prop:prob_imply_Lp}, and to the dominated convergence theorem.




\subsection{The Strong Law of Large Numbers}

\thmpr{Kolmogorov's Zero-One Law}{zero_one_law}{
Let $(X_n)_{n \in \NN}$ be a sequence of independent random variables taking values in arbitrary measurable spaces. For every $n \in \NN$, let $\B_n$ be the $\sigma$-field
\[ \B_n = \sigma(X_k ; k \ge n). \]
Define the asymptotic $\sigma$-field $\B_\infty$ by
\[ \B_\infty = \bigcap_{n=1}^{\infty} \B_n. \]
Then $\PP(B) = 0$ or $1$, for every $B \in \B_\infty$.
}{
For every $n \in \NN$, set
\[ \D_n = \sigma(X_k ; k \le n). \]
Proposition \ref{prop:sigmafield_independence_separation} shows that $\D_n$ is independent of $\B_{n+1}$, hence a fortiori $\D_n$ is independent of $\B_\infty$. Consequently,
\[ \forall A \in \bigcup_{n=1}^{\infty} \D_n,  \quad \forall B \in \B_\infty,  \quad \PP(A \cap B) = \PP(A)\PP(B). \]
Since the class $\bigcup_{n=1}^{\infty} \D_n$ is closed under finite intersections, Proposition \ref{thm:suffisant_condition_sigmafield_independence} implies that $\B_\infty$ is independent of
\[ \sigma\Bigl(\bigcup_{n=1}^{\infty} \D_n\Bigr) = \sigma(X_n ; n \ge 1). \]
In particular, $\B_\infty$ is independent of itself, but this means that, for every $B \in \B_\infty$,
\[\PP(B) = \PP(B \cap B) = \PP(B)^2,\]
which is only possible if $\PP(B) = 0$ or $1$.
}


Let us give a simple illustration of Theorem \ref{thm:zero_one_law}. In the setting of this theorem, a random variable $Y$ taking values in $[-\infty,\infty]$ which is $\B_\infty$-measurable is necessarily equal to a constant almost surely (this can be verified by observing that the distribution function $x \mapsto \PP(Y \le x)$ takes values in $\{0,1\}$). Let us then consider a sequence $(X_n)_{n \in \NN}$ of independent real random variables. The random variable
\[
Z := \limsup_{n \to \infty} \frac{1}{n}(X_1 + \cdots + X_n),
\]
which takes values in $[-\infty,\infty]$, is $\B_\infty$-measurable, since we have also, for every $k \in \NN$,
\[
Z = \limsup_{n \to \infty} \frac{1}{n}(X_k + X_{k+1} + \cdots + X_n),
\]
which shows that $Z$ is $\B_k$-measurable for every $k \in \NN$. Theorem \ref{thm:zero_one_law} thus implies that $Z$ is constant almost surely. In particular, if we know that $\frac{1}{n}(X_1 + \cdots + X_n)$ converges almost surely, the limit must be constant.

\thmpr{Strong Law of Large Numbers}{strong_law_large_numbers}{
Let $(X_n)_{n \ge 1}$ be a sequence of independent and identically distributed real random variables in $L^1$. Then
\[ \frac{1}{n}(X_1 + \cdots + X_n) \conv[a.s.]{n} \EE[X_1]. \]
}{
To simplify notation, we set $S_n = X_1 + \cdots + X_n$ for every $n \in \NN$, and $S_0 = 0$. Let $a > \EE[X_1]$, and
\[ M := \sup_{n \in \ZZ_+} (S_n - na), \]
which is a random variable with values in $[0,\infty]$. \\ 
We will show that
\[ M < \infty, \quad \text{a.s.} \tag{i} \]

Assume that (1) holds (for any choice of $a > \EE[X_1]$). \\ 
Since $S_n \le na + M$ for every $n \in \NN$, it immediately follows from (1) that
\[ \limsup_{n \to \infty} \frac{1}{n} S_n \le a, \quad \text{a.s.} \]
By considering a sequence of values of $a$ that decreases to $\EE[X_1]$, we find
\[ \limsup_{n \to \infty} \frac{1}{n} S_n \le \EE[X_1], \quad \text{a.s.} \]
Replacing $X_n$ by $-X_n$, we get the reverse inequality
\[ \liminf_{n \to \infty} \frac{1}{n} S_n \ge \EE[X_1], \quad \text{a.s.} \]
and the convergence in the theorem follows.

It remains to prove (1). With the notation of Theorem \ref{thm:zero_one_law}, we first observe that $\{M < \infty\}$ belongs to $\B_\infty$. Indeed, we have, for every integer $k \in \NN$,
\[ \{M < \infty\} = \left\{ \sup_{n \in \ZZ_+} (S_n - na) < \infty \right\} = \left\{ \sup_{n \ge k} (S_n - S_k - na) < \infty \right\}, \]
and the event in the right-hand side is $\B_{k+1}$-measurable. It will therefore be enough to prove that $\PP(M < \infty) > 0$, or equivalently that $\PP(M = \infty) < 1$. In the remaining part of the proof we verify that $\PP(M = \infty) < 1$.

For every $k \in \NN$, set
\[ M_k := \sup_{0 \le n \le k} (S_n - na), \qquad M_k' := \sup_{0 \le n \le k} (S_{n+1} - S_1 - na). \]
Then, $M_k$ and $M_k'$ are nonnegative random variables with the same distribution, as a consequence of the fact that the random vectors $(X_1,\ldots,X_k)$ and $(X_2,\ldots,X_{k+1})$ have the same law. It follows that
\[ M = \lim_{k \to \infty} \uparrow\, M_k \qquad \text{and} \qquad M' := \lim_{k \to \infty} \uparrow\, M_k' \]
also have the same law (just observe that $\PP(M' \le x) = \lim \downarrow \PP(M_k' \le x) = \lim \downarrow \PP(M_k \le x) = \PP(M \le x)$ for every $x \in \RR$).

On the other hand, it follows from the definitions that, for every integer $k \ge 0$,
\[ M_{k+1} = \sup\left(0,\ \sup_{1 \le n \le k+1} (S_n - na)\right) = \sup(0, M_k' + X_1 - a), \]
which can also be written as
\[ M_{k+1} = M_k' - \inf(a - X_1, M_k'). \]
Since $M_k'$ has the same law as $M_k$ (and both these variables are clearly in $L^1$), we get
\[ \EE[\inf(a - X_1, M_k')] = \EE[M_k'] - \EE[M_{k+1}] = \EE[M_k] - \EE[M_{k+1}] \le 0 \]
thanks to the trivial inequality $M_k \le M_{k+1}$. We may now apply the dominated convergence theorem to the sequence $\inf(a - X_1, M_k')$, $k \in \NN$, noting that these variables are bounded in absolute value by $|a - X_1|$ (recall that $M_k' \ge 0$). It follows that
\[ \EE[\inf(a - X_1, M')] = \lim_{k \to \infty} \EE[\inf(a - X_1, M_k')] \le 0. \]

We now argue by contradiction to prove that $\PP(M = \infty) < 1$. Suppose that $\PP(M = \infty) = 1$, then we have also $\PP(M' = \infty) = 1$, since $M$ and $M'$ have the same law, and it follows that $\inf(a - X_1, M') = a - X_1$ a.s. But then the inequality in the last display gives $\EE[a - X_1] \le 0$, which is a contradiction since we chose $a > \EE[X_1]$. This contradiction completes the proof.
}

\textbf{Remarks.}
\begin{itemize}
    \item  The $L^1$ integrability assumption is optimal in the sense that it is needed for the limit $\EE[X_1]$ to be defined (and finite). If we remove the $L^1$ assumption but assume instead that the random variables $X_n$ are nonnegative and $\EE[X_1] = \infty$, it is easy to obtain that
    \[ \frac{1}{n}(X_1 + \cdots + X_n) \xrightarrow[n \to \infty]{\mathrm{a.s.}} +\infty, \]
    by applying the theorem to the variables $X_n \wedge K$, for every $K \in \NN$.

    \item  One can show that the convergence of the theorem also holds in $L^1$. The proof will be given at the end of next chapter as an application of martingale theory. As we already mentioned, the almost sure convergence is the most interesting one from a probabilistic point of view.
\end{itemize}



\subsection{Convergence in Distribution}
Recall that $C_b(\RR^d)$ denote the set of all bounded continuous functions from $\RR^d$ into $\RR$. We equip $C_b(\RR^d)$ with the supremum norm
\[ \|\varphi\|=\sup_{x\in\RR^d}|\varphi(x)|. \]

We let $M_1(\RR^d)$ denote the set of all probability measures on $(\RR^d,\B(\RR^d))$.

\defnr{Convergence in distribution}{dist_conv}{A sequence $(\QQ_n)_{n\in\NN}$ in $M_1(\RR^d)$ converges weakly to $\QQ\in M_1(\RR^d)$ if
\[ \forall \varphi\in C_b(\RR^d), \qquad \int \varphi\, d\QQ_n \conv{n} \int \varphi\, d\QQ . \]
A sequence $(X_n)_{n\in\NN}$ of random variables with values in $\RR^d$ converges in distribution to a random variable $X$ with values in $\RR^d$ if the sequence $(P_{X_n})_{n\in\NN}$ converges weakly to $P_X$.\\
By Proposition 8.5, this is equivalent to saying that
\[\forall \varphi\in C_b(\RR^d), \qquad \EE[\varphi(X_n)] \goto{n} \EE[\varphi(X)].\]

We will write $\QQ_n \conv[(w)]{n} \QQ$, resp. $X_n \conv[\D]{n} X$, if the sequence $(\QQ_n)_{n\in\NN}$ converges weakly to $\QQ$, resp. if $(X_n)_{n\in\NN}$ converges in distribution to $X$.
}



\textbf{Remark}
There is some abuse of terminology in saying that the sequence $(X_n)_{n\in\NN}$ converges in distribution to $X$, because the limiting random variable $X$ is not determined uniquely (only its distribution $P_X$ is determined). For this reason, we will sometimes write that a sequence $(X_n)_{n\in\NN}$ of random variables converges in distribution to a probability measure $\QQ$ --- one should of course understand that the laws $P_{X_n}$ converge weakly to $\QQ$. We may also note that the convergence in distribution makes sense even if the random variables $X_n$, $n\in\NN$ are defined on different probability spaces (in the present book, we will however always assume that they are defined on the same probability space). This makes the convergence in distribution very different from the other types of convergence studied in this section.


\proppr{prob_imply_dist}{
    If  
    \[X_n\conv[\PP]{n} X\] 
    Then : 
    \[X_n\conv[\D]{n} X\] 
}{
Suppose first that $X_n$ converges a.s. to $X$. Then, for every $\varphi\in C_b(\RR^d)$, $\varphi(X_n)$ converges a.s. to $\varphi(X)$ and the dominated convergence theorem gives $\EE[\varphi(X_n)] \to \EE[\varphi(X)]$.

If $X_n$ only converges in probability to $X$, then we know from Proposition 10.3 that there is a subsequence of $(X_n)_{n\in\NN}$ that converges a.s. to $X$, and thus, for every $\varphi\in C_b(\RR^d)$, $\EE[\varphi(X_n)]$ converges to $\EE[\varphi(X)]$ along this subsequence. But then, the same argument shows that, from any subsequence of $(X_n)_{n\in\NN}$, we can extract a subsequence along which $\EE[\varphi(X_n)]$ converges to $\EE[\varphi(X)]$. This is only possible if $\EE[\varphi(X_n)] \to \EE[\varphi(X)]$ as $n\to\infty$. }

\textbf{Remark:} The converse to the proposition is (of course) false, because, as we already mentioned, the convergence in distribution of $X_n$ to $X$ does not determine the limiting variable $X$. 


\thmpr{Portmanteau Theorem}{Portmanteau_theorem}{Let $(\QQ_n)_{n\in\NN}$ be a sequence in $M_1(\RR^d)$ and let $\QQ\in M_1(\RR^d)$. The following four assertions are equivalent:
\begin{enumerate}
    \item  The sequence $(\QQ_n)$ converges weakly to $\QQ$.
    \item  For every open subset $G$ of $\RR^d$,
    \[\liminf \QQ_n(G) \ge \QQ(G).\]
    \item  For every closed subset $F$ of $\RR^d$,
    \[\limsup \QQ_n(F) \le \QQ(F).\]
    \item  For every Borel subset $B$ of $\RR^d$ such that $\QQ(\partial B)=0$,
    \[\lim \QQ_n(B) = \QQ(B).\]
\end{enumerate}

}{
Let us start by proving that (1)$\Rightarrow$(2). If $G$ is an open subset of $\RR^d$, we can find a sequence $(\varphi_p)_{p\in\NN}$ of bounded continuous functions such that $0 \le \varphi_p \le \1_G$ and $\varphi_p \uparrow \1_G$ (for instance, $\varphi_p(x)=p\,\mathrm{dist}(x,G^c)\wedge 1$). Then
\[ \liminf_{n\to\infty} \QQ_n(G) \ge \sup_p \left(\liminf_{n\to\infty} \int \varphi_p\, d\QQ_n\right) = \sup_p \left(\int \varphi_p\, d\QQ\right) = \QQ(G). \]

The equivalence (2)$\Leftrightarrow$(3) is immediate since complements of open sets are closed and conversely.

Let us show that (2) and (3) imply (4). If $B\in\B(\RR^d)$, write $\overline{B}$ for the closure of $B$ and $B^\circ$ for the interior of $B$ so that $\partial B=\overline{B}\setminus B^\circ$. By (2) and (3), we have : 
\[\begin{cases}
    \limsup \QQ_n(B) \le& \limsup \QQ_n(\overline{B}) \le \QQ(\overline{B})\\
    \liminf \QQ_n(B) \ge& \liminf \QQ_n(B^\circ) \ge \QQ(B^\circ).
\end{cases}\]

If $\QQ(\partial B)=0$ then $\QQ(\overline{B})=\QQ(B^\circ)=\QQ(B)$ and we get (4).

We still have to prove that (4)$\Rightarrow$(1). Let $\varphi\in C_b(\RR^d)$. Without loss of generality, we can assume that $\varphi \ge 0$. Then, let $K \ge 0$ such that $0 \le \varphi \le K$. By the Fubini theorem,
\[\int \varphi(x)\QQ(dx)= \int \left(\int_0^K \1_{\{t\le\varphi(x)\}} dt \right)\QQ(dx)= \int_0^K \QQ(E_t^\varphi)\,dt,\]
where $E_t^\varphi$ is the closed set $E_t^\varphi := \{x\in\RR^d : \varphi(x)\ge t\}$. Similarly, for every $n$,
\[ \int \varphi(x)\QQ_n(dx) = \int_0^K \QQ_n(E_t^\varphi)\,dt. \]

Now note that $\partial E_t^\varphi \subset \{x\in\RR^d : \varphi(x)=t\}$, and that there are at most countably many values of $t$ such that
\[ \QQ(\{x\in\RR^d : \varphi(x)=t\}) > 0 \]

(indeed, for every $k\in\NN$, there are at most $k$ distinct values of $t$ such that $\QQ(\{x\in\RR^d : \varphi(x)=t\}) \ge \frac{1}{k}$, since the sets $\{x\in\RR^d : \varphi(x)=t\}$ are disjoint when $t$ varies, and $\QQ$ is a probability measure). Hence (4) implies
\[ \QQ_n(E_t^\varphi) \goto{n} \QQ(E_t^\varphi), \]

for every $t\in[0,K]$, except possibly for a countable set of values of $t$, which has zero Lebesgue measure. By dominated convergence, we conclude that
\[ \int \varphi(x)\QQ_n(dx) = \int_0^K \QQ_n(E_t^\varphi)\,dt \goto{n} \int_0^K \QQ(E_t^\varphi)\,dt = \int \varphi(x)\QQ(dx). \]
}


Recall the notation $C_c(\RR^d)$ for the space of all real continuous functions with compact support on $\RR^d$.

\prop{Let $(\QQ_n)_{n\in\NN}$ and $\QQ$ be probability measures on $\RR^d$. Let $H$ be a subset of $C_b(\RR^d)$ whose closure (with respect to the supremum norm) contains $C_c(\RR^d)$. The following properties are equivalent.
\begin{enumerate}
\item  The sequence $(\QQ_n)_{n\in\NN}$ converges weakly to $\QQ$.
\item We have
\[\forall \varphi \in C_c(\RR^d), \qquad \int \varphi\, d\QQ_n \goto{n} \int \varphi\, d\QQ .\]
\item  We have
\[\forall \varphi \in H, \qquad \int \varphi\, d\QQ_n \goto{n} \int \varphi\, d\QQ .\]
\end{enumerate}}{
It is obvious that (1)$\Rightarrow$(2) and (1)$\Rightarrow$(3). Then suppose that (2) holds. Let
$\varphi \in C_b(\RR^d)$ and let $(f_k)_{k\in\NN}$ be a sequence in $C_c(\RR^d)$ such that
$0 \le f_k \le 1$ for every $k \in \NN$, and $f_k \uparrow 1$ as $k \to \infty$. Then, for every $k \in \NN$,
$\varphi f_k \in C_c(\RR^d)$ and thus
\[ \int \varphi f_k \, d\QQ_n \goto{n} \int \varphi f_k \, d\QQ. \tag{i} \]

On the other hand,
\[ \left| \int \varphi \, d\QQ_n - \int \varphi f_k \, d\QQ_n \right| \le \|\varphi\| \int (1-f_k)\, d\QQ_n = \|\varphi\|\left(1-\int f_k\, d\QQ_n\right), \]

and similarly

\begin{align*}
\left| \int \varphi \, d\QQ - \int \varphi f_k \, d\QQ \right|
&\le \|\varphi\|\left(1-\int f_k\, d\QQ\right).
\end{align*}

Hence, using (i), we get for every $k \in \NN$,
\begin{align*}
\limsup_{n\to\infty}
\left| \int \varphi\, d\QQ_n - \int \varphi\, d\QQ \right|
    &\le \limsup_{n\to\infty}
            \left|  \int \varphi\, d\QQ_n - \int \varphi f_k\, d\QQ_n\right|
            + \left| \int \varphi\, d\QQ - \int \varphi f_k\, d\QQ\right| \\
    &\le \|\varphi\|
        \left( \limsup_{n\to\infty} \left(1-\int f_k\, d\QQ_n\right) + \left(1-\int f_k\, d\QQ\right) \right) \\
    &=2\|\varphi\| \left(1-\int f_k\, d\QQ\right).
    \end{align*}

Now we just have to let $k$ tend to $\infty$ (noting that $\int f_k\, d\QQ \uparrow 1$, by monotone convergence) and we find that $\int \varphi\, d\QQ_n$ converges to $\int \varphi\, d\QQ$, so that (1) holds.

We complete the proof by verifying that (3)$\Rightarrow$(2). So assume that (3) holds. If $\varphi \in C_c(\RR^d)$, then, for every $k \in \NN$, we can find a function $\varphi_k \in H$ such that
$\|\varphi-\varphi_k\|\le 1/k$, and it follows that

\begin{align*}
    & \limsup _{n \rightarrow \infty}\left|\int \varphi \mathrm{~d} \QQ_n-\int \varphi \mathrm{d} \QQ\right| \\
    & \leq \limsup _{n \rightarrow \infty}\left(\left|\int\left(\varphi-\varphi_k\right) \mathrm{d} \QQ_n\right|+\left|\int \varphi_k \mathrm{~d} \QQ_n-\int \varphi_k \mathrm{~d} \QQ\right|+\left|\int\left(\varphi_k-\varphi\right) \mathrm{d} \QQ\right|\right) \\
    & \leq \frac{2}{k}
\end{align*}


Since this holds for every $k$, we get that
\[\int \varphi\, d\QQ_n \longrightarrow \int \varphi\, d\QQ.\]
}
Recall our notation $\hat{\QQ}(\xi) = \int e^{i\xi\cdot x}\,\QQ(dx)$ for the Fourier transform of
$\QQ \in M_1(\RR^d)$, and $\Phi_X=\widehat{P_X}$ for the characteristic function of a random variable $X$
with values in $\RR^d$.



\thmpr{Lévy's Theorem}{Levy_theorem}{
A sequence $(\QQ_n)_{n\in\NN}$ in $M_1(\RR^d)$ converges weakly to $\QQ \in M_1(\RR^d)$ if and only if
\[\forall \xi \in \RR^d, \qquad \hat{\QQ}_n(\xi) \xrightarrow[n\to\infty]{} \hat{\QQ}(\xi).
\]

Equivalently, a sequence $(X_n)_{n\in\NN}$ of random variables with values in $\RR^d$ converges in distribution to $X$ if and only if
\[\forall \xi \in \RR^d, \qquad \Phi_{X_n}(\xi) \xrightarrow[n\to\infty]{} \Phi_X(\xi).
\]
}{
It is enough to prove the first assertion. First, if we assume that the sequence
$(\QQ_n)_{n\in\NN}$ converges weakly to $\QQ$, the definition of this convergence shows that
\[ \forall \xi \in \RR^d , \qquad \hat{\QQ}_n(\xi) = \int e^{i\xi\cdot x}\,\QQ_n(dx) \xrightarrow[n\to\infty]{} \int e^{i\xi\cdot x}\,\QQ(dx) = \hat{\QQ}(\xi). \]

Conversely, assume that $\hat{\QQ}_n(\xi)\to \hat{\QQ}(\xi)$ for every $\xi \in \RR^d$ and let us show that
$(\QQ_n)$ converges weakly to $\QQ$. To simplify notation, we only treat the case $d=1$, but
the proof in the general case is exactly similar.

Let $f \in C_c(\RR)$, and, for every $\sigma>0$, let
\[ g_\sigma(x)=\frac{1}{\sigma\sqrt{2\pi}}\exp\!\left(-\frac{x^2}{2\sigma^2}\right). \]

be the density of the $\NN(0,\sigma^2)$ distribution. We already observed in the proof of
Theorem 8.16 that $g_\sigma * f$ converges pointwise to $f$ as $\sigma\to0$. Since $f$ has compact
support here, it is not hard to verify that this convergence is uniform on $\RR$ (use
Proposition 5.8 (i) and note that, for every $\varepsilon>0$, we can find a compact set $K$ such
that $|g_\sigma * f|\le \varepsilon$ for every $x\in\RR\setminus K$ and $\sigma\in(0,1]$).
Thus, if we let $H$ be the subset of $C_b(\RR^d)$ defined by
\[ H:=\{\varphi=g_\sigma * f:\sigma>0 \text{ and } f\in C_c(\RR^d)\} \]

the closure of $H$ in $C_b(\RR^d)$ contains $C_c(\RR^d)$.

On the other hand, we also saw in the proof of Theorem 8.16 that, for any $\nu\in M_1(\RR)$,
\begin{align*}
\int g_\sigma * f \mathrm{~d} v & =\int f(x) g_\sigma * v(x) \mathrm{d} x \\
& =\int f(x)\left((\sigma \sqrt{2 \pi})^{-1} \int e^{i \xi x} g_{1 / \sigma}(\xi) \widehat{v}(-\xi) \mathrm{d} \xi\right) \mathrm{d} x
\end{align*}

We apply this formula to $\nu=\QQ_n$ and $\nu=\QQ$. From our assumption
$\hat{\QQ}_n(\xi)\to\hat{\QQ}(\xi)$ for every $\xi\in\RR$ and to the trivial bound
$|\hat{\QQ}_n(\xi)|\le1$, we can apply the dominated convergence theorem to get that, for every
$x\in\RR$,
\[ \int e^{i\xi x} g_{1/\sigma}(\xi)\hat{\QQ}_n(-\xi)\,d\xi \xrightarrow[n\to\infty]{} \int e^{i\xi x} g_{1/\sigma}(\xi)\hat{\QQ}(-\xi)\,d\xi . \]

Since the quantities in the last display are bounded by $1$ and $f\in C_c(\RR)$, we can use (10.3) and again the dominated convergence theorem to get
\[ \int g_\sigma * f\, d\QQ_n \xrightarrow[n\to\infty]{} \int g_\sigma * f\, d\QQ . \]

Finally, we have proved that $\int \varphi\, d\QQ_n \to \int \varphi\, d\QQ$ for every
$\varphi\in H$, and we know that the closure of $H$ in $C_b(\RR^d)$ contains
$C_c(\RR^d)$. By Proposition 10.12, this is enough to show that the sequence
$(\QQ_n)_{n\in\NN}$ converges weakly to $\QQ$. 
}



\subsection{The Central Limit Theorem}
Let $(X_n)_{n \in \NN}$ be a sequence of independent and identically distributed real random variables in $L^1$. The strong law of large numbers asserts that
\[ \frac{1}{n}(X_1 + \cdots + X_n) \xrightarrow[n \to \infty]{a.s.} \EE[X_1]. \]

We are now interested in the speed of this convergence, meaning that we want information about the typical size of
\[
\frac{1}{n}(X_1 + \cdots + X_n) - \EE[X_1]
\]
when $n$ is large.

Under the additional assumption that the variables $X_n$ are in $L^2$, it is easy to guess the answer, because a simple calculation made in the proof of the weak law of large numbers (Theorem 9.13) gives
\begin{align*}
\EE\left[(X_1 + \cdots + X_n - n\EE[X_1])^2\right]
&= \mathrm{var}(X_1 + \cdots + X_n) \\
&= n\,\mathrm{var}(X_1).
\end{align*}

This shows that the mean value of $(X_1 + \cdots + X_n - n\EE[X_1])^2$ grows linearly with $n$, and thus the typical size of $X_1 + \cdots + X_n - n\EE[X_1]$ is expected to be $\sqrt{n}$, or equivalently the typical size of
\[ \frac{1}{n}(X_1 + \cdots + X_n) - \EE[X_1] \]
should be $1/\sqrt{n}$. The central limit theorem gives a much more precise statement.

\thmpr{Central Limit Theorem}{Theorem 10.15}{
Let $(X_n)_{n \in \NN}$ be a sequence of independent and identically distributed real random variables in $L^2$. Set $\sigma^2 = \mathrm{var}(X_1)$ and assume that $\sigma > 0$. Then
\[ \frac{1}{\sqrt{n}}(X_1 + \cdots + X_n - n\EE[X_1]) \xrightarrow[n \to \infty]{(d)} \NN(0,\sigma^2). \]

where $\NN(0,\sigma^2)$ is the Gaussian distribution with mean $0$ and variance $\sigma^2$. Hence, for every $a,b \in \RR$ with $a < b$,
\[ \lim_{n \to \infty} \PP(X_1 + \cdots + X_n \in [n\EE[X_1] + a\sqrt{n},\, n\EE[X_1] + b\sqrt{n}]) = \frac{1}{\sigma\sqrt{2\pi}} \int_a^b \exp\left(-\frac{x^2}{2\sigma^2}\right)\,dx. \]
}{
The second part of the statement follows from the first one, thanks to Proposition 10.11. To prove the first assertion, note that we can assume $\EE[X_1] = 0$ (just replace $X_n$ by $X_n - \EE[X_n]$). Then set
\[ Z_n = \frac{1}{\sqrt{n}}(X_1 + \cdots + X_n). \]

The characteristic function of $Z_n$ is
\begin{align*}
    \Phi_{Z_n}(\xi) &= \EE\left[\exp\left(i\xi \frac{X_1 + \cdots + X_n}{\sqrt{n}}\right)\right] \\
    &= \EE\left[\exp\left(i\frac{\xi}{\sqrt{n}} X_1\right)\right]^n \\
    &= \Phi_{X_1}\left(\frac{\xi}{\sqrt{n}}\right)^n.
\end{align*}

where, in the second equality, we used the fact that the random variables $X_i$ are independent and identically distributed. On the other hand, by Proposition 8.17, we have
\[ \Phi_{X_1}(\xi) = 1 + i\xi \EE[X_1] - \frac{1}{2}\xi^2 \EE[X_1^2] + o(\xi^2) = 1 - \frac{\sigma^2 \xi^2}{2} + o(\xi^2) \]
as $\xi \to 0$. For every fixed $\xi \in \RR$, we have thus
\[ \Phi_{X_1}\left(\frac{\xi}{\sqrt{n}}\right) = 1 - \frac{\sigma^2 \xi^2}{2n} + o\left(\frac{1}{n}\right) \]
as $n \to \infty$. Thanks to proposition \ref{prop:prob_imply_Lp} and the last display, we have, for every $\xi \in \RR$,
\[ \lim_{n \to \infty} \Phi_{Z_n}(\xi) = \lim_{n \to \infty} \left(1 - \frac{\sigma^2 \xi^2}{2n} + o\left(\frac{1}{n}\right)\right)^n = \exp\left(-\frac{\sigma^2 \xi^2}{2}\right), \]
and, by Lemma 8.15, the limit is the characteristic function of the Gaussian $\NN(0,\sigma^2)$ distribution. An application of Lévy's theorem \ref{thm:Levy_theorem} now completes the proof.
}
\textbf{Remark} When $\sigma = 0$ there is not much to say since the variables $X_n$ are all equal a.s. to the constant $\EE[X_1]$.


\subsection{The Multidimensional Central Limit Theorem}
We now assume that $(X_n)_{n \in \NN}$ is a sequence of independent and identically distributed random variables with values in $\RR^d$, such that $\EE[\|X_1\|] < \infty$. We can apply the strong law of large numbers to each component of $X_n$ and we have
\[
\frac{1}{n}(X_1 + \cdots + X_n) \xrightarrow[n \to \infty]{a.s.} \EE[X_1],
\]
as in the real case. If we assume furthermore that $\EE[\|X_1\|^2] < \infty$, then we expect to have also an analog of the central limit theorem. However, in contrast with the law of large numbers, this is not immediate because the convergence in distribution of a sequence of random vectors does not follow from the convergence in distribution of each component sequence (in fact, the law of the limiting vector is not determined by its marginal distributions, as we already noticed in Chapter 8).

To extend the central limit theorem to the case of random variables with values in $\RR^d$, we must first generalize the Gaussian distribution in higher dimensions.

\textbf{Definition 10.16} A random variable $X$ with values in $\RR^d$ is called a Gaussian vector if any linear combination of the components of $X$ is a (real) Gaussian variable. This is equivalent to the existence of a vector $m \in \RR^d$ and a $d \times d$ symmetric nonnegative definite matrix $C$ such that, for every $\xi \in \RR^d$,
\[
\EE[\exp(i \xi \cdot X)] = \exp\left(i \xi \cdot m - \frac{1}{2}\,{}^t\xi C \xi \right).
\tag{10.5}
\]

Moreover, we have then $\EE[X] = m$ and $K_X = C$, and we say that $X$ follows the Gaussian $\NN_d(m,C)$ distribution.

To establish the equivalence of the two forms of the definition, first assume that any linear combination of the components of $X$ is Gaussian. In particular, these components are in $L^2$ and thus the expectation $\EE[X]$ and the covariance matrix $K_X$ are well-defined. Then, for every $\xi \in \RR^d$, we have $\EE[\xi \cdot X] = \xi \cdot \EE[X]$ and $\mathrm{var}(\xi \cdot X) = {}^t\xi K_X \xi$ (by (8.3)), so that $\xi \cdot X$ has the Gaussian $\NN(\xi \cdot \EE[X], {}^t\xi K_X \xi)$ distribution. Formula (10.5), with $m = \EE[X]$ and $C = K_X$, now follows from Lemma 8.15 giving the characteristic function of a real Gaussian variable. Conversely, if (10.5) holds, then by replacing $\xi$ with $\lambda \xi$ for $\lambda \in \RR$, and using again Lemma 8.15, we get that the characteristic function of $\xi \cdot X$ coincides with the characteristic function of the Gaussian $\NN(\xi \cdot m, {}^t\xi C \xi)$ distribution, so that, in particular, $\xi \cdot X$ is Gaussian.

\textbf{Remark} In order for $X$ to be a Gaussian vector, it is not sufficient that each component of $X$ is a Gaussian (real) random variable. Let us give a simple example with $d = 2$. Let $X_1$ be Gaussian $\NN(0,1)$ and let $U$ be a random variable with law $\PP(U = 1) = \PP(U = -1) = 1/2$, which is independent of $X_1$. If we set $X_2 = U X_1$, then one immediately checks that $X_2$ is also Gaussian $\NN(0,1)$. On the other hand, $(X_1,X_2)$ is certainly not a Gaussian vector, because $\PP(X_1 + X_2 = 0) = \PP(U = -1) = 1/2$, which makes it impossible for $X_1 + X_2$ to be Gaussian.

Let $Y_1,\ldots,Y_d$ be independent Gaussian (real) random variables. Then $(Y_1,\ldots,Y_d)$ is a Gaussian vector. Indeed we noticed at the end of Section 9.5 that any linear combination of independent Gaussian random variables is Gaussian. Alternatively, we can also notice that formula (10.5) holds, with a diagonal matrix $C$. Conversely, if the covariance matrix of a Gaussian vector is diagonal, then its components are independent. This is an immediate application of Proposition 9.8 (iv), and we will come back to this result in the next chapter.

\proppr{idk}{Let $m \in \RR^d$ and let $C$ be a symmetric nonnegative definite $d \times d$ matrix. One can construct a Gaussian $\NN_d(m,C)$ random vector.}{
It is enough to treat the case $m = 0$ (if $X$ is Gaussian $\NN_d(0,C)$, $X + m$ is Gaussian $\NN_d(m,C)$). Set $A = \sqrt{C}$ in such a way that $A$ is a symmetric nonnegative definite matrix and $A^2 = C$. Let $Y_1,\ldots,Y_d$ be independent Gaussian $\NN(0,1)$ variables, and $Y = (Y_1,\ldots,Y_d)$, so that the covariance matrix $K_Y$ is the identity matrix. Then $X = AY$ follows the Gaussian $\NN_d(0,C)$ distribution. Indeed any linear combination of the components of $X$ is also a linear combination of $Y_1,\ldots,Y_d$ and is therefore Gaussian. Moreover, by (8.2),
\[
K_X = A K_Y {}^t A = A^2 = C,
\]
which shows that $X$ is Gaussian $\NN_d(0,C)$.}

\thmpr{Multidimensional Central Limit Theorem}{Theorem 10.18}{
Let $(X_n)_{n \in \NN}$ be a sequence of independent and identically distributed random variables with values in $\RR^d$, such that $\EE[\|X_1\|^2] < \infty$. Then,
\[
\frac{1}{\sqrt{n}}(X_1 + \cdots + X_n - n\EE[X_1])
\xrightarrow[n \to \infty]{(d)} \NN_d(0, K_{X_1}).
\]
}{
Replacing $X_n$ by $X_n - \EE[X_n]$ allows us to assume that $\EE[X_1] = 0$. For every $\xi \in \RR^d$, the central limit theorem in the real case (Theorem 10.15) shows that
\[
\frac{1}{\sqrt{n}}(\xi \cdot X_1 + \cdots + \xi \cdot X_n)
\xrightarrow[n \to \infty]{(d)} \NN(0,\sigma^2),
\]
where $\sigma^2 = \EE[(\xi \cdot X_1)^2] = {}^t\xi K_{X_1} \xi$. This implies that
\[
\EE\left[\exp\left(i\xi \cdot \frac{X_1 + \cdots + X_n}{\sqrt{n}}\right)\right]
\longrightarrow \exp\left(-\frac{1}{2}\,{}^t\xi K_{X_1}\xi\right)
\]
and Lévy's theorem (Theorem 10.13) gives the desired result.
}




\end{document}